% Options for packages loaded elsewhere
% Options for packages loaded elsewhere
\PassOptionsToPackage{unicode}{hyperref}
\PassOptionsToPackage{hyphens}{url}
\PassOptionsToPackage{dvipsnames,svgnames,x11names}{xcolor}
%
\documentclass[
  ngerman,
  a4paper,
  DIV=11]{scrreprt}
\usepackage{xcolor}
\usepackage[margin=2.5cm]{geometry}
\usepackage{amsmath,amssymb}
\setcounter{secnumdepth}{5}
\usepackage{iftex}
\ifPDFTeX
  \usepackage[T1]{fontenc}
  \usepackage[utf8]{inputenc}
  \usepackage{textcomp} % provide euro and other symbols
\else % if luatex or xetex
  \usepackage{unicode-math} % this also loads fontspec
  \defaultfontfeatures{Scale=MatchLowercase}
  \defaultfontfeatures[\rmfamily]{Ligatures=TeX,Scale=1}
\fi
\usepackage{lmodern}
\ifPDFTeX\else
  % xetex/luatex font selection
\fi
% Use upquote if available, for straight quotes in verbatim environments
\IfFileExists{upquote.sty}{\usepackage{upquote}}{}
\IfFileExists{microtype.sty}{% use microtype if available
  \usepackage[]{microtype}
  \UseMicrotypeSet[protrusion]{basicmath} % disable protrusion for tt fonts
}{}
\makeatletter
\@ifundefined{KOMAClassName}{% if non-KOMA class
  \IfFileExists{parskip.sty}{%
    \usepackage{parskip}
  }{% else
    \setlength{\parindent}{0pt}
    \setlength{\parskip}{6pt plus 2pt minus 1pt}}
}{% if KOMA class
  \KOMAoptions{parskip=half}}
\makeatother
% Make \paragraph and \subparagraph free-standing
\makeatletter
\ifx\paragraph\undefined\else
  \let\oldparagraph\paragraph
  \renewcommand{\paragraph}{
    \@ifstar
      \xxxParagraphStar
      \xxxParagraphNoStar
  }
  \newcommand{\xxxParagraphStar}[1]{\oldparagraph*{#1}\mbox{}}
  \newcommand{\xxxParagraphNoStar}[1]{\oldparagraph{#1}\mbox{}}
\fi
\ifx\subparagraph\undefined\else
  \let\oldsubparagraph\subparagraph
  \renewcommand{\subparagraph}{
    \@ifstar
      \xxxSubParagraphStar
      \xxxSubParagraphNoStar
  }
  \newcommand{\xxxSubParagraphStar}[1]{\oldsubparagraph*{#1}\mbox{}}
  \newcommand{\xxxSubParagraphNoStar}[1]{\oldsubparagraph{#1}\mbox{}}
\fi
\makeatother

\usepackage{color}
\usepackage{fancyvrb}
\newcommand{\VerbBar}{|}
\newcommand{\VERB}{\Verb[commandchars=\\\{\}]}
\DefineVerbatimEnvironment{Highlighting}{Verbatim}{commandchars=\\\{\}}
% Add ',fontsize=\small' for more characters per line
\usepackage{framed}
\definecolor{shadecolor}{RGB}{241,243,245}
\newenvironment{Shaded}{\begin{snugshade}}{\end{snugshade}}
\newcommand{\AlertTok}[1]{\textcolor[rgb]{0.68,0.00,0.00}{#1}}
\newcommand{\AnnotationTok}[1]{\textcolor[rgb]{0.37,0.37,0.37}{#1}}
\newcommand{\AttributeTok}[1]{\textcolor[rgb]{0.40,0.45,0.13}{#1}}
\newcommand{\BaseNTok}[1]{\textcolor[rgb]{0.68,0.00,0.00}{#1}}
\newcommand{\BuiltInTok}[1]{\textcolor[rgb]{0.00,0.23,0.31}{#1}}
\newcommand{\CharTok}[1]{\textcolor[rgb]{0.13,0.47,0.30}{#1}}
\newcommand{\CommentTok}[1]{\textcolor[rgb]{0.37,0.37,0.37}{#1}}
\newcommand{\CommentVarTok}[1]{\textcolor[rgb]{0.37,0.37,0.37}{\textit{#1}}}
\newcommand{\ConstantTok}[1]{\textcolor[rgb]{0.56,0.35,0.01}{#1}}
\newcommand{\ControlFlowTok}[1]{\textcolor[rgb]{0.00,0.23,0.31}{\textbf{#1}}}
\newcommand{\DataTypeTok}[1]{\textcolor[rgb]{0.68,0.00,0.00}{#1}}
\newcommand{\DecValTok}[1]{\textcolor[rgb]{0.68,0.00,0.00}{#1}}
\newcommand{\DocumentationTok}[1]{\textcolor[rgb]{0.37,0.37,0.37}{\textit{#1}}}
\newcommand{\ErrorTok}[1]{\textcolor[rgb]{0.68,0.00,0.00}{#1}}
\newcommand{\ExtensionTok}[1]{\textcolor[rgb]{0.00,0.23,0.31}{#1}}
\newcommand{\FloatTok}[1]{\textcolor[rgb]{0.68,0.00,0.00}{#1}}
\newcommand{\FunctionTok}[1]{\textcolor[rgb]{0.28,0.35,0.67}{#1}}
\newcommand{\ImportTok}[1]{\textcolor[rgb]{0.00,0.46,0.62}{#1}}
\newcommand{\InformationTok}[1]{\textcolor[rgb]{0.37,0.37,0.37}{#1}}
\newcommand{\KeywordTok}[1]{\textcolor[rgb]{0.00,0.23,0.31}{\textbf{#1}}}
\newcommand{\NormalTok}[1]{\textcolor[rgb]{0.00,0.23,0.31}{#1}}
\newcommand{\OperatorTok}[1]{\textcolor[rgb]{0.37,0.37,0.37}{#1}}
\newcommand{\OtherTok}[1]{\textcolor[rgb]{0.00,0.23,0.31}{#1}}
\newcommand{\PreprocessorTok}[1]{\textcolor[rgb]{0.68,0.00,0.00}{#1}}
\newcommand{\RegionMarkerTok}[1]{\textcolor[rgb]{0.00,0.23,0.31}{#1}}
\newcommand{\SpecialCharTok}[1]{\textcolor[rgb]{0.37,0.37,0.37}{#1}}
\newcommand{\SpecialStringTok}[1]{\textcolor[rgb]{0.13,0.47,0.30}{#1}}
\newcommand{\StringTok}[1]{\textcolor[rgb]{0.13,0.47,0.30}{#1}}
\newcommand{\VariableTok}[1]{\textcolor[rgb]{0.07,0.07,0.07}{#1}}
\newcommand{\VerbatimStringTok}[1]{\textcolor[rgb]{0.13,0.47,0.30}{#1}}
\newcommand{\WarningTok}[1]{\textcolor[rgb]{0.37,0.37,0.37}{\textit{#1}}}

\usepackage{longtable,booktabs,array}
\usepackage{calc} % for calculating minipage widths
% Correct order of tables after \paragraph or \subparagraph
\usepackage{etoolbox}
\makeatletter
\patchcmd\longtable{\par}{\if@noskipsec\mbox{}\fi\par}{}{}
\makeatother
% Allow footnotes in longtable head/foot
\IfFileExists{footnotehyper.sty}{\usepackage{footnotehyper}}{\usepackage{footnote}}
\makesavenoteenv{longtable}
\usepackage{graphicx}
\makeatletter
\newsavebox\pandoc@box
\newcommand*\pandocbounded[1]{% scales image to fit in text height/width
  \sbox\pandoc@box{#1}%
  \Gscale@div\@tempa{\textheight}{\dimexpr\ht\pandoc@box+\dp\pandoc@box\relax}%
  \Gscale@div\@tempb{\linewidth}{\wd\pandoc@box}%
  \ifdim\@tempb\p@<\@tempa\p@\let\@tempa\@tempb\fi% select the smaller of both
  \ifdim\@tempa\p@<\p@\scalebox{\@tempa}{\usebox\pandoc@box}%
  \else\usebox{\pandoc@box}%
  \fi%
}
% Set default figure placement to htbp
\def\fps@figure{htbp}
\makeatother


% definitions for citeproc citations
\NewDocumentCommand\citeproctext{}{}
\NewDocumentCommand\citeproc{mm}{%
  \begingroup\def\citeproctext{#2}\cite{#1}\endgroup}
\makeatletter
 % allow citations to break across lines
 \let\@cite@ofmt\@firstofone
 % avoid brackets around text for \cite:
 \def\@biblabel#1{}
 \def\@cite#1#2{{#1\if@tempswa , #2\fi}}
\makeatother
\newlength{\cslhangindent}
\setlength{\cslhangindent}{1.5em}
\newlength{\csllabelwidth}
\setlength{\csllabelwidth}{3em}
\newenvironment{CSLReferences}[2] % #1 hanging-indent, #2 entry-spacing
 {\begin{list}{}{%
  \setlength{\itemindent}{0pt}
  \setlength{\leftmargin}{0pt}
  \setlength{\parsep}{0pt}
  % turn on hanging indent if param 1 is 1
  \ifodd #1
   \setlength{\leftmargin}{\cslhangindent}
   \setlength{\itemindent}{-1\cslhangindent}
  \fi
  % set entry spacing
  \setlength{\itemsep}{#2\baselineskip}}}
 {\end{list}}
\usepackage{calc}
\newcommand{\CSLBlock}[1]{\hfill\break\parbox[t]{\linewidth}{\strut\ignorespaces#1\strut}}
\newcommand{\CSLLeftMargin}[1]{\parbox[t]{\csllabelwidth}{\strut#1\strut}}
\newcommand{\CSLRightInline}[1]{\parbox[t]{\linewidth - \csllabelwidth}{\strut#1\strut}}
\newcommand{\CSLIndent}[1]{\hspace{\cslhangindent}#1}

\ifLuaTeX
\usepackage[bidi=basic]{babel}
\else
\usepackage[bidi=default]{babel}
\fi
% get rid of language-specific shorthands (see #6817):
\let\LanguageShortHands\languageshorthands
\def\languageshorthands#1{}
\ifLuaTeX
  \usepackage[german]{selnolig} % disable illegal ligatures
\fi


\setlength{\emergencystretch}{3em} % prevent overfull lines

\providecommand{\tightlist}{%
  \setlength{\itemsep}{0pt}\setlength{\parskip}{0pt}}



 


\KOMAoption{captions}{tableheading}
\makeatletter
\@ifpackageloaded{tcolorbox}{}{\usepackage[skins,breakable]{tcolorbox}}
\@ifpackageloaded{fontawesome5}{}{\usepackage{fontawesome5}}
\definecolor{quarto-callout-color}{HTML}{909090}
\definecolor{quarto-callout-note-color}{HTML}{0758E5}
\definecolor{quarto-callout-important-color}{HTML}{CC1914}
\definecolor{quarto-callout-warning-color}{HTML}{EB9113}
\definecolor{quarto-callout-tip-color}{HTML}{00A047}
\definecolor{quarto-callout-caution-color}{HTML}{FC5300}
\definecolor{quarto-callout-color-frame}{HTML}{acacac}
\definecolor{quarto-callout-note-color-frame}{HTML}{4582ec}
\definecolor{quarto-callout-important-color-frame}{HTML}{d9534f}
\definecolor{quarto-callout-warning-color-frame}{HTML}{f0ad4e}
\definecolor{quarto-callout-tip-color-frame}{HTML}{02b875}
\definecolor{quarto-callout-caution-color-frame}{HTML}{fd7e14}
\makeatother
\makeatletter
\@ifpackageloaded{bookmark}{}{\usepackage{bookmark}}
\makeatother
\makeatletter
\@ifpackageloaded{caption}{}{\usepackage{caption}}
\AtBeginDocument{%
\ifdefined\contentsname
  \renewcommand*\contentsname{Inhaltsverzeichnis}
\else
  \newcommand\contentsname{Inhaltsverzeichnis}
\fi
\ifdefined\listfigurename
  \renewcommand*\listfigurename{Abbildungsverzeichnis}
\else
  \newcommand\listfigurename{Abbildungsverzeichnis}
\fi
\ifdefined\listtablename
  \renewcommand*\listtablename{Tabellenverzeichnis}
\else
  \newcommand\listtablename{Tabellenverzeichnis}
\fi
\ifdefined\figurename
  \renewcommand*\figurename{Abbildung}
\else
  \newcommand\figurename{Abbildung}
\fi
\ifdefined\tablename
  \renewcommand*\tablename{Tabelle}
\else
  \newcommand\tablename{Tabelle}
\fi
}
\@ifpackageloaded{float}{}{\usepackage{float}}
\floatstyle{ruled}
\@ifundefined{c@chapter}{\newfloat{codelisting}{h}{lop}}{\newfloat{codelisting}{h}{lop}[chapter]}
\floatname{codelisting}{Listing}
\newcommand*\listoflistings{\listof{codelisting}{Listingverzeichnis}}
\makeatother
\makeatletter
\makeatother
\makeatletter
\@ifpackageloaded{caption}{}{\usepackage{caption}}
\@ifpackageloaded{subcaption}{}{\usepackage{subcaption}}
\makeatother
\usepackage{bookmark}
\IfFileExists{xurl.sty}{\usepackage{xurl}}{} % add URL line breaks if available
\urlstyle{same}
\hypersetup{
  pdftitle={Interaktive Skripte mit KI erstellen},
  pdfauthor={Prof.~Dr.~Roman Bartnik},
  pdflang={de},
  colorlinks=true,
  linkcolor={blue},
  filecolor={Maroon},
  citecolor={Blue},
  urlcolor={Blue},
  pdfcreator={LaTeX via pandoc}}


\title{Interaktive Skripte mit KI erstellen}
\usepackage{etoolbox}
\makeatletter
\providecommand{\subtitle}[1]{% add subtitle to \maketitle
  \apptocmd{\@title}{\par {\large #1 \par}}{}{}
}
\makeatother
\subtitle{Werkstatt-Skript zum Workshop am 11.05.2026, TH Köln}
\author{Prof.~Dr.~Roman Bartnik}
\date{11.05.2026}
\begin{document}
\maketitle

\renewcommand*\contentsname{Inhaltsverzeichnis}
{
\hypersetup{linkcolor=}
\setcounter{tocdepth}{2}
\tableofcontents
}

\bookmarksetup{startatroot}

\chapter{Willkommen}\label{willkommen}

Werkstatt-Skript zum Workshop am 11.05.2026

\hfill\break

\bookmarksetup{startatroot}

\chapter{Was du heute baust}\label{was-du-heute-baust}

Am Ende des Tages steht ein eigenes interaktives Skript live im Netz. Du
hast es selbst gerendert, gepusht und veröffentlicht. Das ist das Ziel
--- alles andere ist Mittel zum Zweck.

\subsubsection{Am Ende des Tages kannst du
\ldots{}}\label{am-ende-des-tages-kannst-du}

\begin{enumerate}
\def\labelenumi{\arabic{enumi}.}
\tightlist
\item
  ein Quarto-Projekt lokal in \textbf{RStudio} anlegen, rendern und über
  \textbf{GitHub Pages} veröffentlichen,
\item
  ein \textbf{automatisches Inhaltsverzeichnis} erzeugen und Quellen aus
  \textbf{Zotero} zitieren,
\item
  das Skript im \textbf{TH-Köln-Stil} formatieren (\texttt{\_brand.yml},
  CSS),
\item
  ein \textbf{interaktives Diagramm} und ein \textbf{iFrame-Spiel}
  einbauen,
\item
  einen \textbf{KI-Tutor} verlinken, der die Lernenden im Skript
  begleitet.
\end{enumerate}

\section{Drei Referenzen}\label{drei-referenzen}

Bevor wir bauen, schauen wir, was möglich ist. Drei Skripte zeigen die
Bandbreite --- vom Kapitel zum Buch zum Kurs.

\subsection{Lakens --- Statistical
Inferences}\label{lakens-statistical-inferences}

Ein vollständiges Open-Access-Lehrbuch zur Statistik mit Quizzen,
Code-Demos und Videos. Maßstab für „so gut kann ein Skript werden''.

\href{https://lakens.github.io/statistical_inferences/}{Zum Skript →}

\subsection{Bartnik --- KI als Hilfe zum Lehren und
Lernen}\label{bartnik-ki-als-hilfe-zum-lehren-und-lernen}

Mein eigenes Skript, mit interaktiven Widgets, eingebetteten Tutor-Bots
und H5P-Quizzen.

\href{https://th-koln-bartnik.github.io/genai4teaching/}{Zum Skript →}

\subsection{Békés --- Data Analysis with
AI}\label{buxe9kuxe9s-data-analysis-with-ai}

Ein vollständiger Kurs samt Folien, Notebooks und Übungen. Maßstab für
„Skript plus Slides plus Code''.

\href{https://gabors-data-analysis.com/ai-course/}{Zum Skript →}

\section{Wähle deinen
Beispielstrang}\label{wuxe4hle-deinen-beispielstrang}

Damit du nicht auch noch Inhalte erfinden musst, nutzen wir ein
durchgehendes Mini-Thema: eine \textbf{kurze Einführung ins
wissenschaftliche Arbeiten}. Wähle einen der drei Stränge und behalte
ihn den ganzen Tag bei.

\begin{longtable}[]{@{}
  >{\raggedright\arraybackslash}p{(\linewidth - 4\tabcolsep) * \real{0.3333}}
  >{\raggedright\arraybackslash}p{(\linewidth - 4\tabcolsep) * \real{0.3333}}
  >{\raggedright\arraybackslash}p{(\linewidth - 4\tabcolsep) * \real{0.3333}}@{}}
\toprule\noalign{}
\begin{minipage}[b]{\linewidth}\raggedright
Strang
\end{minipage} & \begin{minipage}[b]{\linewidth}\raggedright
Worum es geht
\end{minipage} & \begin{minipage}[b]{\linewidth}\raggedright
Quelle
\end{minipage} \\
\midrule\noalign{}
\endhead
\bottomrule\noalign{}
\endlastfoot
\textbf{A --- Wissenschaftliche Einstellung} & Skepsis, Belege,
intellektuelle Demut &
\href{https://publish.obsidian.md/wagp-calling-bullshit/wissenschaftliche+Einstellung}{Calling
Bullshit, Kap. „Wissenschaftliche Einstellung''} \\
\textbf{B --- Recherche und Quellenqualität} & Suchen, prüfen, Quellen
einordnen &
\href{https://publish.obsidian.md/wagp-calling-bullshit/Kritische+Informationssuche}{Calling
Bullshit, Kap. „Kritische Informationssuche''} \\
\textbf{C --- Transparent schreiben} & Methoden offenlegen,
Reproduzierbarkeit &
\href{https://publish.obsidian.md/wagp-calling-bullshit/Transparent+schreiben}{Calling
Bullshit, Kap. „Transparent schreiben''} \\
\end{longtable}

Die Wahl ist Geschmackssache. Die technischen Schritte sind in jedem
Strang identisch.

\section{Tagesablauf}\label{tagesablauf}

\begin{longtable}[]{@{}
  >{\raggedright\arraybackslash}p{(\linewidth - 4\tabcolsep) * \real{0.3333}}
  >{\raggedright\arraybackslash}p{(\linewidth - 4\tabcolsep) * \real{0.3333}}
  >{\raggedright\arraybackslash}p{(\linewidth - 4\tabcolsep) * \real{0.3333}}@{}}
\toprule\noalign{}
\begin{minipage}[b]{\linewidth}\raggedright
Uhrzeit
\end{minipage} & \begin{minipage}[b]{\linewidth}\raggedright
Block
\end{minipage} & \begin{minipage}[b]{\linewidth}\raggedright
Was passiert
\end{minipage} \\
\midrule\noalign{}
\endhead
\bottomrule\noalign{}
\endlastfoot
09:00--09:45 & Auftakt & Vorstellung, Vorerfahrung, Ziele, Tour durch
die drei Vorbilder, Strangwahl \\
09:45--10:30 & Kap. 1 --- Didaktik & Drei Lerntheorie-Argumente,
Lernspiel zur Selbstprüfung \\
10:30--11:45 & Kap. 2 --- Setup & Bausteine, Memory-Spiel, lokale
Grundstruktur, \textbf{erster Push live} \\
11:45--12:30 & Pause & \\
12:30--13:15 & Kap. 3a --- Übersicht & Inhaltsverzeichnis,
Zotero-Bibliothek, automatische Zitate \\
13:15--14:00 & Kap. 3b --- Stil & TH-Köln-Brand einbauen \\
14:00--14:45 & Kap. 4 --- Charts & Interaktives Diagramm mit echten
Daten \\
14:45--15:30 & Kap. 5 --- iFrames + Tutor & H5P-/iFrame-Einbettung, Link
zum KI-Tutor \\
15:30--16:00 & Showcase + Puffer & Wer früh durch ist, baut Slides (Kap.
6); alle anderen holen auf \\
\end{longtable}

\section{Wie du dieses Skript
benutzt}\label{wie-du-dieses-skript-benutzt}

Jedes Kapitel öffnet mit einem Lernziel-Kasten („Was du nach diesem
Kapitel kannst''), gefolgt von einem \textbf{„Mach mit''}-Block --- das
ist die eigentliche Übung. Dazwischen steht so wenig Text wie möglich.
Am Ende jedes Kapitels findest du einen Trouble-Shooting-Kasten („Was,
wenn \ldots``) und einen fertigen \textbf{Tutor-Prompt} zum Kopieren.

\begin{tcolorbox}[enhanced jigsaw, opacitybacktitle=0.6, leftrule=.75mm, title=\textcolor{quarto-callout-tip-color}{\faLightbulb}\hspace{0.5em}{RStudio ist deine Arbeitszentrale}, colframe=quarto-callout-tip-color-frame, arc=.35mm, toprule=.15mm, colbacktitle=quarto-callout-tip-color!10!white, bottomrule=.15mm, colback=white, rightrule=.15mm, breakable, toptitle=1mm, bottomtitle=1mm, opacityback=0, titlerule=0mm, coltitle=black, left=2mm]

Innerhalb von RStudio hast du \textbf{Editor}, \textbf{Terminal} (für
\texttt{quarto\ render} und \texttt{quarto\ preview}),
\textbf{Files-Pane} (zum Anlegen von Ordnern und Dateien),
\textbf{Build-Pane} (Knopf „Render Book'') und \textbf{Console} (für
R-Pakete) auf einem Schirm. Daneben brauchst du nur \textbf{GitHub
Desktop} zum Versionieren und einen Browser-Tab mit dem \textbf{Tutor}.

\href{https://chatgpt.com/g/g-69fec93348108191af1ffb8ce0cf6712-tutor-interaktive-skripte}{\textbf{Tutor
öffnen --- „Interaktive Skripte''} →}

\end{tcolorbox}

\section{Ordnerstruktur deines
Skripts}\label{ordnerstruktur-deines-skripts}

Damit du immer weißt, wo etwas hingehört (Lakens-Konvention):

\begin{longtable}[]{@{}
  >{\raggedright\arraybackslash}p{(\linewidth - 2\tabcolsep) * \real{0.2297}}
  >{\raggedright\arraybackslash}p{(\linewidth - 2\tabcolsep) * \real{0.7703}}@{}}
\toprule\noalign{}
\begin{minipage}[b]{\linewidth}\raggedright
Ordner
\end{minipage} & \begin{minipage}[b]{\linewidth}\raggedright
Was kommt rein?
\end{minipage} \\
\midrule\noalign{}
\endhead
\bottomrule\noalign{}
\endlastfoot
\texttt{parts/} & Deine Kapiteldateien (\texttt{.qmd}) \\
\texttt{images/} & Logos, Screenshots, statische Grafiken (PNG, JPG,
SVG) \\
\texttt{interactions/} & HTML-Widgets ohne Backend (Lernspiel, Memory,
Quiz) \\
\texttt{include/} & Wiederverwendbare qmd-Schnipsel (per
\texttt{\{\{\textless{}\ include\ \textgreater{}\}\}}) \\
\texttt{data/} & CSV/JSON-Daten für interaktive Charts \\
\texttt{appendix/} & Quickstart, Troubleshooting, Tutor-Prompt \\
\end{longtable}

\section{Voraussetzungen}\label{voraussetzungen}

Wenn du noch nichts installiert hast, springe zuerst in den
\href{appendix/quickstart.qmd}{Quick-Starter Guide}. Dort steht, was du
in welcher Reihenfolge installierst (Quarto, R, RStudio, GitHub Desktop,
Zotero) --- mit Screenshots.

\section{Quellen}\label{quellen}

Lakens (2022); Bartnik (2025); Békés (2025); Biggs \& Tang (2011); Mayer
(2021).

\part{Teil 1 --- Werkstatt}

\chapter{1 --- Warum interaktiv?}\label{warum-interaktiv}

\subsubsection{Was du nach diesem Kapitel
kannst}\label{was-du-nach-diesem-kapitel-kannst}

\begin{itemize}
\tightlist
\item
  drei lerntheoretisch gestützte Argumente nennen, \textbf{warum} ein
  interaktives Skript ein passives Skriptum schlägt,
\item
  jedes Argument einem \textbf{konkreten Skript-Element} zuordnen (Quiz,
  Diagramm, Tutor),
\item
  entscheiden, \textbf{wann sich Interaktivität lohnt} und wann nicht.
\end{itemize}

\section{Drei Argumente in drei
Sätzen}\label{drei-argumente-in-drei-suxe4tzen}

Lehrforschung sagt es seit Jahren --- aber selten so kompakt: Lernen
klappt besser, wenn Lernende \textbf{selbst etwas tun},
\textbf{regelmäßig abrufen} und \textbf{Bilder neben Worten} bekommen.
Drei Prinzipien, drei Belege, drei Konsequenzen für dein Skript.

\subsection{Generatives Lernen}\label{generatives-lernen}

Wer den Stoff selbst zusammenfasst, erklärt oder anwendet, behält ihn
besser, als wer ihn nur liest (\citeproc{ref-fiorella2016}{Fiorella \&
Mayer, 2016}). Das ist kein Geschmack, das ist eine Meta-Analyse über
acht Lernaktivitäten --- Zusammenfassen, Selbsterklären, Lehren, Mappen,
Vorhersagen, Vorstellen, Zeichnen, Aufführen. Übersetzt: in deinem
Skript braucht jedes Kapitel \textbf{eine Aufgabe, die der Lernende
selbst tut}, nicht nur einen Text, den er liest.

\subsection{Retrieval Practice}\label{retrieval-practice}

Sich an etwas erinnern zu \textbf{müssen} verbessert das Behalten
stärker als wiederholtes Lesen (\citeproc{ref-roediger2006}{Roediger \&
Karpicke, 2006}). Karpicke und Roediger ließen Studierende dasselbe
Material drei Mal lesen oder zwei Mal lesen plus einmal aus dem
Gedächtnis aufschreiben --- die zweite Gruppe schlug die erste eine
Woche später deutlich. Übersetzt: ein \textbf{kurzes Quiz} oder eine
\textbf{Drag-and-Drop-Aufgabe} in deinem Skript wirkt wie eine
Mini-Klausur, ohne dass jemand die Klausur korrigieren muss.

\subsection{Multimedia-Prinzip}\label{multimedia-prinzip}

Worte plus passende Bilder schlagen Worte allein
(\citeproc{ref-mayer2021}{Mayer, 2021}). „Passend'' heißt: das Bild
zeigt, wovon der Text spricht --- nicht zur Dekoration, sondern als
zweite Repräsentation desselben Inhalts. Übersetzt: das
\textbf{interaktive Diagramm}, das du in Kapitel 4 einbauen wirst, hat
einen lerntheoretischen Grund --- keine Spielerei.

\section{Was das für dein Skript
heißt}\label{was-das-fuxfcr-dein-skript-heiuxdft}

Drei Bauteile, die du heute alle einbaust, lassen sich genau diesen drei
Prinzipien zuordnen:

\begin{longtable}[]{@{}
  >{\raggedright\arraybackslash}p{(\linewidth - 4\tabcolsep) * \real{0.3333}}
  >{\raggedright\arraybackslash}p{(\linewidth - 4\tabcolsep) * \real{0.3333}}
  >{\raggedright\arraybackslash}p{(\linewidth - 4\tabcolsep) * \real{0.3333}}@{}}
\toprule\noalign{}
\begin{minipage}[b]{\linewidth}\raggedright
Prinzip
\end{minipage} & \begin{minipage}[b]{\linewidth}\raggedright
Skript-Element
\end{minipage} & \begin{minipage}[b]{\linewidth}\raggedright
Wo im Workshop?
\end{minipage} \\
\midrule\noalign{}
\endhead
\bottomrule\noalign{}
\endlastfoot
Generatives Lernen & KI-Tutor, eigene Übungsaufgabe & Kap. 5 \\
Retrieval Practice & Lernspiel, Memory, Quiz (iFrame oder H5P) & Kap. 1,
2, 5 \\
Multimedia & Interaktives Diagramm, Grafiken & Kap. 4 \\
\end{longtable}

Das ist die didaktische Architektur deines Skripts --- und gleichzeitig
die Architektur dieses Workshops.

\section{\texorpdfstring{Wann \textbf{kein} interaktives
Element?}{Wann kein interaktives Element?}}\label{wann-kein-interaktives-element}

Drei Situationen, in denen Interaktivität schadet:

\begin{itemize}
\tightlist
\item
  \textbf{Kognitive Überforderung.} Wenn die Lernenden gleichzeitig
  neuen Stoff verarbeiten \emph{und} die Bedienung des Widgets lernen
  müssen, kostet das Aufmerksamkeit, die für den Inhalt fehlt
  (\citeproc{ref-mayer2021}{Mayer, 2021}). Lieber erst Stoff, dann
  Übung.
\item
  \textbf{Reine Illustration.} Wenn das Diagramm keine Variation
  erlaubt, die didaktisch interessant ist, reicht ein PNG.
  Interaktivität sollte einen Lernzweck haben, nicht nur „können wir
  technisch''.
\item
  \textbf{Ablenkung durch Kontrolle.} Wenn ein Spiel die Lernenden in
  den Spielmodus zieht („wie viele Punkte krieg ich?{}``), kann das vom
  Inhalt wegführen. Faustregel: Belohnung für \textbf{richtige Inhalte},
  nicht für Schnelligkeit.
\end{itemize}

\section{Mach mit --- Lernspiel zur
Selbstprüfung}\label{mach-mit-lernspiel-zur-selbstpruxfcfung}

Drei Beispiele, drei Prinzipien, fünf Minuten. Ziehe jedes Beispiel auf
das Prinzip, das es illustriert.

\begin{tcolorbox}[enhanced jigsaw, opacitybacktitle=0.6, leftrule=.75mm, title=\textcolor{quarto-callout-tip-color}{\faLightbulb}\hspace{0.5em}{Lernspiel öffnen}, colframe=quarto-callout-tip-color-frame, arc=.35mm, toprule=.15mm, colbacktitle=quarto-callout-tip-color!10!white, bottomrule=.15mm, colback=white, rightrule=.15mm, breakable, toptitle=1mm, bottomtitle=1mm, opacityback=0, titlerule=0mm, coltitle=black, left=2mm]

\href{../interactions/lernspiel-kapitel1.html}{Lernspiel: Welches
Prinzip steckt hinter welchem Beispiel? →}

\end{tcolorbox}

\begin{tcolorbox}[enhanced jigsaw, opacitybacktitle=0.6, leftrule=.75mm, title=\textcolor{quarto-callout-warning-color}{\faExclamationTriangle}\hspace{0.5em}{Was, wenn \ldots{}}, colframe=quarto-callout-warning-color-frame, arc=.35mm, toprule=.15mm, colbacktitle=quarto-callout-warning-color!10!white, bottomrule=.15mm, colback=white, rightrule=.15mm, breakable, toptitle=1mm, bottomtitle=1mm, opacityback=0, titlerule=0mm, coltitle=black, left=2mm]

\begin{itemize}
\tightlist
\item
  \textbf{Das Spiel öffnet sich nicht?} → Direkt im Browser über die
  Datei \texttt{interactions/lernspiel-kapitel1.html} öffnen. Die Datei
  ist eine einzige HTML-Datei, ohne Backend.
\item
  \textbf{Die Lösungen scheinen mehrdeutig?} → Stimmt --- in der Praxis
  überschneiden sich die Prinzipien. Schau in die Erklärung am Ende des
  Spiels, dort steht, warum die jeweilige Zuordnung am stärksten passt.
\end{itemize}

\end{tcolorbox}

\section{Tutor}\label{tutor}

\begin{tcolorbox}[enhanced jigsaw, opacitybacktitle=0.6, leftrule=.75mm, title=\textcolor{quarto-callout-tip-color}{\faLightbulb}\hspace{0.5em}{Tutor öffnen}, colframe=quarto-callout-tip-color-frame, arc=.35mm, toprule=.15mm, colbacktitle=quarto-callout-tip-color!10!white, bottomrule=.15mm, colback=white, rightrule=.15mm, breakable, toptitle=1mm, bottomtitle=1mm, opacityback=0, titlerule=0mm, coltitle=black, left=2mm]

\href{https://chatgpt.com/g/g-69fec93348108191af1ffb8ce0cf6712-tutor-interaktive-skripte}{Tutor
„Interaktive Skripte'' →}

Frag den Tutor nach passenden interaktiven Elementen für dein Kapitel.
Vorschlag-Prompt:

\begin{quote}
Ich überlege, welche interaktiven Elemente in mein Skript zum Thema
{[}Strang A/B/C{]} passen. Frag mich nach dem Lernziel des jeweiligen
Kapitels und schlage dann pro Kapitel EIN passendes Element vor ---
verbunden mit einem der drei Lernprinzipien (generativ, retrieval,
multimedia). Begründe in einem Satz, warum genau dieses Element passt.
\end{quote}

\end{tcolorbox}

\chapter*{Literatur}\label{literatur}
\addcontentsline{toc}{chapter}{Literatur}

\markboth{Literatur}{Literatur}

\phantomsection\label{refs}
\begin{CSLReferences}{1}{0}
\bibitem[\citeproctext]{ref-allaire2024}
Allaire, J. J., Teague, C., Scheidegger, C., Xie, Y., \& Dervieux, C.
(2024). \emph{Quarto: An Open-Source Scientific and Technical Publishing
System}. \url{https://quarto.org/}

\bibitem[\citeproctext]{ref-bergstrom2024}
Bergstrom, C. T., \& West, J. D. (2024). \emph{Calling Bullshit:
Wissenschaftliches Arbeiten und kritisches Denken}.
\url{https://publish.obsidian.md/wagp-calling-bullshit/}

\bibitem[\citeproctext]{ref-bryan2018}
Bryan, J. (2018). Excuse Me, Do You Have a Moment to Talk About Version
Control? \emph{The American Statistician}, \emph{72}(1), 20--27.
\url{https://doi.org/10.1080/00031305.2017.1399928}

\bibitem[\citeproctext]{ref-fiorella2016}
Fiorella, L., \& Mayer, R. E. (2016). Eight Ways to Promote Generative
Learning. \emph{Educational Psychology Review}, \emph{28}(4), 717--741.
\url{https://doi.org/10.1007/s10648-015-9348-9}

\bibitem[\citeproctext]{ref-healy2018}
Healy, K. (2018). \emph{Data Visualization: A Practical Introduction}.
Princeton University Press.

\bibitem[\citeproctext]{ref-karpicke2011}
Karpicke, J. D., \& Blunt, J. R. (2011). Retrieval Practice Produces
More Learning than Elaborative Studying with Concept Mapping.
\emph{Science}, \emph{331}(6018), 772--775.
\url{https://doi.org/10.1126/science.1199327}

\bibitem[\citeproctext]{ref-kasneci2023}
Kasneci, E., Sessler, K., Küchemann, S., Bannert, M., Dementieva, D.,
Fischer, F., Gasser, U., Groh, G., Günnemann, S., Hüllermeier, E.,
Krusche, S., Kutyniok, G., Michaeli, T., Nerdel, C., Pfeffer, J.,
Poquet, O., Sailer, M., Schmidt, A., Seidel, T., \ldots{} Kasneci, G.
(2023). {ChatGPT} for Good? On Opportunities and Challenges of Large
Language Models for Education. \emph{Learning and Individual
Differences}, \emph{103}, 102274.
\url{https://doi.org/10.1016/j.lindif.2023.102274}

\bibitem[\citeproctext]{ref-lakens2022}
Lakens, D. (2022). \emph{Improving Your Statistical Inferences}.
\url{https://doi.org/10.5281/zenodo.6409077}

\bibitem[\citeproctext]{ref-mayer2021}
Mayer, R. E. (2021). \emph{Multimedia Learning} (3. Aufl.). Cambridge
University Press. \url{https://doi.org/10.1017/9781316941355}

\bibitem[\citeproctext]{ref-mollick2023}
Mollick, E. R., \& Mollick, L. (2023). Assigning {AI}: Seven Approaches
for Students, with Prompts. \emph{SSRN Working Paper}.
\url{https://doi.org/10.2139/ssrn.4475995}

\bibitem[\citeproctext]{ref-roediger2006}
Roediger, H. L., \& Karpicke, J. D. (2006). Test-Enhanced Learning:
Taking Memory Tests Improves Long-Term Retention. \emph{Psychological
Science}, \emph{17}(3), 249--255.
\url{https://doi.org/10.1111/j.1467-9280.2006.01693.x}

\bibitem[\citeproctext]{ref-owid2024}
Roser, M., Ritchie, H., Mathieu, E., \& Ortiz-Ospina, E. (2024).
\emph{Our World in Data}. \url{https://ourworldindata.org/}

\bibitem[\citeproctext]{ref-tufte2001}
Tufte, E. R. (2001). \emph{The Visual Display of Quantitative
Information} (2. Aufl.). Graphics Press.

\bibitem[\citeproctext]{ref-wickham2019}
Wickham, H., Averick, M., Bryan, J., Chang, W., McGowan, L. D.,
François, R., Grolemund, G., Hayes, A., Henry, L., Hester, J., Kuhn, M.,
Pedersen, T. L., Miller, E., Bache, S. M., Müller, K., Ooms, J.,
Robinson, D., Seidel, D. P., Spinu, V., \ldots{} Yutani, H. (2019).
Welcome to the Tidyverse. \emph{Journal of Open Source Software},
\emph{4}(43), 1686. \url{https://doi.org/10.21105/joss.01686}

\end{CSLReferences}

\chapter{2 --- Bausteine und Workflow}\label{bausteine-und-workflow}

\subsubsection{Was du nach diesem Kapitel
kannst}\label{was-du-nach-diesem-kapitel-kannst-1}

\begin{itemize}
\tightlist
\item
  die sechs Bausteine \textbf{Quarto, YAML, JSON, HTML, iFrame, GitHub}
  in eigenen Worten erklären,
\item
  ein leeres Quarto-Projekt \textbf{lokal anlegen} (RStudio + Terminal),
\item
  es mit \textbf{GitHub Desktop} in ein neues Repository legen,
\item
  es einmal über \textbf{GitHub Pages} veröffentlichen --- ein erster
  klickbarer Link.
\end{itemize}

\section{Worum es im Kern geht}\label{worum-es-im-kern-geht}

Du schreibst Text in einer Datei. Ein Programm verwandelt diese Datei in
eine Webseite. Eine zweite Datei legt fest, wie sie aussieht. Ein
Online-Speicher liefert die Webseite weltweit aus. Mehr ist es nicht ---
alles, was wir heute machen, sind Variationen dieses
Vier-Zeilen-Workflows. Dass Forschungsarbeit von Versionsverwaltung
profitiert, ist seit langem belegt (\citeproc{ref-bryan2018}{Bryan,
2018}); Quarto bringt diese Logik in die Lehre
(\citeproc{ref-allaire2024}{Allaire et al., 2024}).

\section{Die sechs Bausteine --- Lay-Erklärung
zuerst}\label{die-sechs-bausteine-lay-erkluxe4rung-zuerst}

Jeder Begriff wird in der Fachsprache schwerer als nötig. Hier zuerst
die einfache Idee, dann das Detail.

\begin{longtable}[]{@{}
  >{\raggedright\arraybackslash}p{(\linewidth - 4\tabcolsep) * \real{0.3333}}
  >{\raggedright\arraybackslash}p{(\linewidth - 4\tabcolsep) * \real{0.3333}}
  >{\raggedright\arraybackslash}p{(\linewidth - 4\tabcolsep) * \real{0.3333}}@{}}
\toprule\noalign{}
\begin{minipage}[b]{\linewidth}\raggedright
Baustein
\end{minipage} & \begin{minipage}[b]{\linewidth}\raggedright
In einem Satz
\end{minipage} & \begin{minipage}[b]{\linewidth}\raggedright
Was es technisch ist
\end{minipage} \\
\midrule\noalign{}
\endhead
\bottomrule\noalign{}
\endlastfoot
\textbf{Quarto} & Programm, das aus einer Textdatei eine Webseite, ein
PDF oder Folien macht. & Open-Source-Publishing-System, Nachfolger von R
Markdown (\citeproc{ref-allaire2024}{Allaire et al., 2024}). \\
\textbf{YAML} & Einstellungsdatei am Anfang jedes Kapitels --- sagt z.
B., ob es ein Inhaltsverzeichnis gibt. & „YAML Ain't Markup Language'',
strukturiertes Format mit Einrückungen. \\
\textbf{JSON} & Wie YAML, aber mit Klammern. Wird intern für Daten
verwendet. & „JavaScript Object Notation'', Standardformat zum
Datenaustausch. \\
\textbf{HTML} & Sprache, in der Webseiten geschrieben sind. & „HyperText
Markup Language'', Tags wie \texttt{\textless{}h1\textgreater{}} oder
\texttt{\textless{}p\textgreater{}}. \\
\textbf{iFrame} & Fenster im Fenster --- bettet eine fremde Seite in
deine ein. & HTML-Tag
\texttt{\textless{}iframe\ src="..."\textgreater{}}, isoliert vom Rest
der Seite. \\
\textbf{GitHub} & Speicher für Code im Netz, mit eingebautem Webserver
(Pages). & Hosting-Plattform auf Basis von Git, gehört Microsoft. \\
\end{longtable}

\begin{tcolorbox}[enhanced jigsaw, opacitybacktitle=0.6, leftrule=.75mm, title=\textcolor{quarto-callout-tip-color}{\faLightbulb}\hspace{0.5em}{Mach mit --- Memory-Spiel zu den Bausteinen}, colframe=quarto-callout-tip-color-frame, arc=.35mm, toprule=.15mm, colbacktitle=quarto-callout-tip-color!10!white, bottomrule=.15mm, colback=white, rightrule=.15mm, breakable, toptitle=1mm, bottomtitle=1mm, opacityback=0, titlerule=0mm, coltitle=black, left=2mm]

Drei Minuten, sechs Paare. Versuche, alle Begriffe ihrer Lay-Definition
zuzuordnen.

\href{../interactions/memory-spiel.html}{Memory-Spiel öffnen →}

\end{tcolorbox}

\section{Der Workflow in einem Bild}\label{der-workflow-in-einem-bild}

Drei Werkzeuge, ein Kreislauf:

\begin{verbatim}
RStudio (schreiben + rendern)  →  GitHub Desktop (versionieren)  →  GitHub Pages (veröffentlichen)
        ↑                                                                            │
        └──────────────────────────  Tutor im Browser-Tab (helfen)  ─────────────────┘
\end{verbatim}

Versionsverwaltung ist hier kein Nice-to-have. Sie ist die Voraussetzung
dafür, dass dein Skript jederzeit reproduzierbar gerendert werden kann
und dass du ohne Angst Änderungen vornimmst --- ein Punkt, den Bryan
(\citeproc{ref-bryan2018}{2018}) für die Forschungsarbeit ausführlich
begründet, der für die Lehre genauso gilt.

\section{RStudio als Arbeitszentrale --- vier
Bereiche}\label{rstudio-als-arbeitszentrale-vier-bereiche}

Du arbeitest fast immer in RStudio. Vier Bereiche reichen für alles, was
wir heute machen:

\begin{longtable}[]{@{}
  >{\raggedright\arraybackslash}p{(\linewidth - 4\tabcolsep) * \real{0.3333}}
  >{\raggedright\arraybackslash}p{(\linewidth - 4\tabcolsep) * \real{0.3333}}
  >{\raggedright\arraybackslash}p{(\linewidth - 4\tabcolsep) * \real{0.3333}}@{}}
\toprule\noalign{}
\begin{minipage}[b]{\linewidth}\raggedright
Bereich
\end{minipage} & \begin{minipage}[b]{\linewidth}\raggedright
Wo?
\end{minipage} & \begin{minipage}[b]{\linewidth}\raggedright
Wofür?
\end{minipage} \\
\midrule\noalign{}
\endhead
\bottomrule\noalign{}
\endlastfoot
\textbf{Editor} & oben links & \texttt{.qmd}-Dateien schreiben \\
\textbf{Terminal} & unten links, Tab neben Console &
\texttt{quarto\ render}, \texttt{quarto\ preview} \\
\textbf{Files} & unten rechts & neue Ordner und Dateien anlegen, Dateien
öffnen \\
\textbf{Build} & oben rechts & Knopf „Render Book'' als
Klick-Alternative \\
\end{longtable}

Drei nützliche Shortcuts:

\begin{itemize}
\tightlist
\item
  \textbf{Strg+Shift+K} --- die aktuell geöffnete \texttt{.qmd}-Datei
  einzeln rendern (für schnelles Vorschauen einzelner Kapitel).
\item
  \textbf{Strg+Shift+B} --- komplettes Buch rendern (entspricht dem
  Build-Knopf).
\item
  \textbf{Alt+Shift+R} --- neuen Terminal-Tab öffnen, falls keiner
  sichtbar ist.
\end{itemize}

\begin{tcolorbox}[enhanced jigsaw, opacitybacktitle=0.6, leftrule=.75mm, title=\textcolor{quarto-callout-note-color}{\faInfo}\hspace{0.5em}{Hinweis}, colframe=quarto-callout-note-color-frame, arc=.35mm, toprule=.15mm, colbacktitle=quarto-callout-note-color!10!white, bottomrule=.15mm, colback=white, rightrule=.15mm, breakable, toptitle=1mm, bottomtitle=1mm, opacityback=0, titlerule=0mm, coltitle=black, left=2mm]

\textbf{Falls dein Terminal-Tab fehlt:} \textbf{Tools → Terminal → New
Terminal}. Voreinstellung auf Windows ist PowerShell, einstellbar unter
\textbf{Tools → Global Options → Terminal}. Quarto-Befehle funktionieren
in jeder Shell.

\end{tcolorbox}

\section{Mach mit --- die ersten 30
Minuten}\label{mach-mit-die-ersten-30-minuten}

Halte dich an diese Reihenfolge. Bei Problemen: Tutor fragen oder zur
Trouble-Shooting-Box am Ende springen.

\subsection{Schritt 1 --- neues Projekt in
RStudio}\label{schritt-1-neues-projekt-in-rstudio}

\begin{enumerate}
\def\labelenumi{\arabic{enumi}.}
\tightlist
\item
  RStudio öffnen → \textbf{File → New Project → New Directory → Quarto
  Book}.
\item
  Ordnernamen wählen, z. B. \texttt{mein-erstes-skript}.
\item
  Speicherort wählen --- am besten direkt in deinem GitHub-Ordner.
\item
  \textbf{Create Project}.
\end{enumerate}

RStudio legt automatisch eine Mini-Buchstruktur an
(\texttt{\_quarto.yml}, \texttt{index.qmd}, ein paar Beispielkapitel).

\subsection{Schritt 2 --- Render und Preview (in
RStudio)}\label{schritt-2-render-und-preview-in-rstudio}

Im \textbf{Terminal-Tab} (unten neben Console):

\begin{Shaded}
\begin{Highlighting}[]
\ExtensionTok{quarto}\NormalTok{ preview}
\end{Highlighting}
\end{Shaded}

Im Browser öffnet sich \texttt{http://localhost:4XXX} mit einer
Live-Vorschau. Jede Änderung an einer \texttt{.qmd}-Datei führt zu
sofortigem Reload. Alternative ohne Tippen: \textbf{Build-Pane} oben
rechts → Knopf „Render Book'' (rendert das ganze Buch einmal, ohne
Live-Reload).

\subsection{Schritt 3 --- Beispielkapitel auf deinen Strang
anpassen}\label{schritt-3-beispielkapitel-auf-deinen-strang-anpassen}

Öffne \texttt{intro.qmd} und ersetze den Beispielinhalt durch eine
Mini-Einleitung zu deinem gewählten Strang. Drei Vorlagen, je nach
Strangwahl:

\subsubsection{Strang A --- Wissenschaftliche Einstellung}

\begin{Shaded}
\begin{Highlighting}[]
\FunctionTok{\# Wissenschaftliche Einstellung}

\NormalTok{Wissenschaftliche Einstellung heißt, Aussagen nicht zu glauben, sondern}
\NormalTok{zu prüfen — und sich selbst genauso zu prüfen wie andere. Drei Haltungen}
\NormalTok{gehören dazu: Skepsis, intellektuelle Demut, Bereitschaft zur Korrektur.}
\NormalTok{Quelle: Bergstrom \& West (2024).}
\end{Highlighting}
\end{Shaded}

\subsubsection{Strang B --- Recherche und Quellenqualität}

\begin{Shaded}
\begin{Highlighting}[]
\FunctionTok{\# Recherche und Quellenqualität}

\NormalTok{Quellen unterscheiden sich nicht nur darin, ob sie wahr oder falsch sind,}
\NormalTok{sondern darin, wie nachprüfbar sie sind. In diesem Kapitel: drei Filter,}
\NormalTok{mit denen du eine Quelle in 30 Sekunden grob einordnest.}
\NormalTok{Quelle: Bergstrom \& West (2024).}
\end{Highlighting}
\end{Shaded}

\subsubsection{Strang C --- Transparent schreiben}

\begin{Shaded}
\begin{Highlighting}[]
\FunctionTok{\# Transparent schreiben}

\NormalTok{Transparenz heißt, dass jeder Leser nachvollziehen kann, wie du zu deinen}
\NormalTok{Ergebnissen gekommen bist — welche Daten, welche Methoden, welche}
\NormalTok{Annahmen. In diesem Kapitel: drei Bausteine eines transparenten}
\NormalTok{Methodenteils. Quelle: Bergstrom \& West (2024).}
\end{Highlighting}
\end{Shaded}

\subsection{Schritt 4 --- Ordnerstruktur anlegen (in
RStudio)}\label{schritt-4-ordnerstruktur-anlegen-in-rstudio}

Im \textbf{Files-Pane} (unten rechts) klickst du auf \textbf{New Folder}
und legst nacheinander an: \texttt{images/}, \texttt{interactions/},
\texttt{include/}, \texttt{data/}. Das sind die Ordner für deine
späteren Bilder, HTML-Widgets, eingebundene Schnipsel und Datendateien
--- die Aufteilung folgt der Konvention, die Lakens
(\citeproc{ref-lakens2022}{2022}) in seinem Skript verwendet.

\subsection{Schritt 5 --- GitHub Desktop
verknüpfen}\label{schritt-5-github-desktop-verknuxfcpfen}

\begin{enumerate}
\def\labelenumi{\arabic{enumi}.}
\tightlist
\item
  GitHub Desktop öffnen → \textbf{File → Add Local Repository} → Ordner
  wählen.
\item
  Wenn noch kein Git-Repo: „Create a repository'' akzeptieren.
\item
  \textbf{Publish repository} klicken → Name eingeben (z. B.
  \texttt{mein-erstes-skript}) → \textbf{Public} → \textbf{Publish}.
\end{enumerate}

Dein Code liegt jetzt auf GitHub.

\subsection{Schritt 6 --- GitHub Pages
aktivieren}\label{schritt-6-github-pages-aktivieren}

\begin{enumerate}
\def\labelenumi{\arabic{enumi}.}
\tightlist
\item
  Auf github.com im Repository auf \textbf{Settings → Pages}.
\item
  \textbf{Source: Deploy from a branch}.
\item
  \textbf{Branch: main → Folder: /docs → Save}.
\item
  Nach 1--2 Minuten: dein Skript liegt unter
  \texttt{https://DEIN-USERNAME.github.io/mein-erstes-skript/}.
\end{enumerate}

\begin{tcolorbox}[enhanced jigsaw, opacitybacktitle=0.6, leftrule=.75mm, title=\textcolor{quarto-callout-warning-color}{\faExclamationTriangle}\hspace{0.5em}{Was, wenn \ldots{}}, colframe=quarto-callout-warning-color-frame, arc=.35mm, toprule=.15mm, colbacktitle=quarto-callout-warning-color!10!white, bottomrule=.15mm, colback=white, rightrule=.15mm, breakable, toptitle=1mm, bottomtitle=1mm, opacityback=0, titlerule=0mm, coltitle=black, left=2mm]

\begin{itemize}
\tightlist
\item
  \textbf{Render funktioniert lokal, aber GitHub Pages zeigt 404?} → In
  \texttt{\_quarto.yml} muss \texttt{output-dir:\ docs} stehen, und der
  Ordner \texttt{docs/} muss committet sein. Schau in
  \texttt{.gitignore}, ob \texttt{/docs/} ausgeschlossen wird, und nimm
  die Zeile heraus.
\item
  \textbf{YAML-Fehler beim Rendern?} → Einrückungen prüfen (zwei
  Leerzeichen, keine Tabs). Doppelpunkte und Anführungszeichen
  ebenfalls.
\item
  \textbf{GitHub Desktop fragt nach Login?} → Im Browser anmelden, dann
  in GitHub Desktop „Sign in''.
\item
  \textbf{\texttt{quarto:\ command\ not\ found}?} → Quarto-Installation
  prüfen mit \texttt{quarto\ -\/-version}. Wenn leer: neu installieren
  von \url{https://quarto.org/docs/get-started/}.
\end{itemize}

\end{tcolorbox}

\section{Tutor}\label{tutor-1}

\begin{tcolorbox}[enhanced jigsaw, opacitybacktitle=0.6, leftrule=.75mm, title=\textcolor{quarto-callout-tip-color}{\faLightbulb}\hspace{0.5em}{Tutor öffnen}, colframe=quarto-callout-tip-color-frame, arc=.35mm, toprule=.15mm, colbacktitle=quarto-callout-tip-color!10!white, bottomrule=.15mm, colback=white, rightrule=.15mm, breakable, toptitle=1mm, bottomtitle=1mm, opacityback=0, titlerule=0mm, coltitle=black, left=2mm]

\href{https://chatgpt.com/g/g-69fec93348108191af1ffb8ce0cf6712-tutor-interaktive-skripte}{Tutor
„Interaktive Skripte'' →}

Vorschlag-Prompt für dieses Kapitel:

\begin{quote}
Ich versuche gerade, ein Quarto-Buch lokal in RStudio anzulegen und über
GitHub Pages zu veröffentlichen. Ich bin Anfänger, kein IT-Experte.
Erkläre mir Schritt für Schritt, was ich als Nächstes tun muss, immer
mit einer Lay-Erklärung VOR dem Code-Schnipsel. Mein aktueller Stand:
{[}hier kurz beschreiben{]}. Mein Ziel: einen klickbaren Link auf mein
Skript. Frag mich zurück, wenn etwas unklar ist.
\end{quote}

\end{tcolorbox}

\chapter*{Literatur}\label{literatur-1}
\addcontentsline{toc}{chapter}{Literatur}

\markboth{Literatur}{Literatur}

\phantomsection\label{refs}
\begin{CSLReferences}{1}{0}
\bibitem[\citeproctext]{ref-allaire2024}
Allaire, J. J., Teague, C., Scheidegger, C., Xie, Y., \& Dervieux, C.
(2024). \emph{Quarto: An Open-Source Scientific and Technical Publishing
System}. \url{https://quarto.org/}

\bibitem[\citeproctext]{ref-bergstrom2024}
Bergstrom, C. T., \& West, J. D. (2024). \emph{Calling Bullshit:
Wissenschaftliches Arbeiten und kritisches Denken}.
\url{https://publish.obsidian.md/wagp-calling-bullshit/}

\bibitem[\citeproctext]{ref-bryan2018}
Bryan, J. (2018). Excuse Me, Do You Have a Moment to Talk About Version
Control? \emph{The American Statistician}, \emph{72}(1), 20--27.
\url{https://doi.org/10.1080/00031305.2017.1399928}

\bibitem[\citeproctext]{ref-fiorella2016}
Fiorella, L., \& Mayer, R. E. (2016). Eight Ways to Promote Generative
Learning. \emph{Educational Psychology Review}, \emph{28}(4), 717--741.
\url{https://doi.org/10.1007/s10648-015-9348-9}

\bibitem[\citeproctext]{ref-healy2018}
Healy, K. (2018). \emph{Data Visualization: A Practical Introduction}.
Princeton University Press.

\bibitem[\citeproctext]{ref-karpicke2011}
Karpicke, J. D., \& Blunt, J. R. (2011). Retrieval Practice Produces
More Learning than Elaborative Studying with Concept Mapping.
\emph{Science}, \emph{331}(6018), 772--775.
\url{https://doi.org/10.1126/science.1199327}

\bibitem[\citeproctext]{ref-kasneci2023}
Kasneci, E., Sessler, K., Küchemann, S., Bannert, M., Dementieva, D.,
Fischer, F., Gasser, U., Groh, G., Günnemann, S., Hüllermeier, E.,
Krusche, S., Kutyniok, G., Michaeli, T., Nerdel, C., Pfeffer, J.,
Poquet, O., Sailer, M., Schmidt, A., Seidel, T., \ldots{} Kasneci, G.
(2023). {ChatGPT} for Good? On Opportunities and Challenges of Large
Language Models for Education. \emph{Learning and Individual
Differences}, \emph{103}, 102274.
\url{https://doi.org/10.1016/j.lindif.2023.102274}

\bibitem[\citeproctext]{ref-lakens2022}
Lakens, D. (2022). \emph{Improving Your Statistical Inferences}.
\url{https://doi.org/10.5281/zenodo.6409077}

\bibitem[\citeproctext]{ref-mayer2021}
Mayer, R. E. (2021). \emph{Multimedia Learning} (3. Aufl.). Cambridge
University Press. \url{https://doi.org/10.1017/9781316941355}

\bibitem[\citeproctext]{ref-mollick2023}
Mollick, E. R., \& Mollick, L. (2023). Assigning {AI}: Seven Approaches
for Students, with Prompts. \emph{SSRN Working Paper}.
\url{https://doi.org/10.2139/ssrn.4475995}

\bibitem[\citeproctext]{ref-roediger2006}
Roediger, H. L., \& Karpicke, J. D. (2006). Test-Enhanced Learning:
Taking Memory Tests Improves Long-Term Retention. \emph{Psychological
Science}, \emph{17}(3), 249--255.
\url{https://doi.org/10.1111/j.1467-9280.2006.01693.x}

\bibitem[\citeproctext]{ref-owid2024}
Roser, M., Ritchie, H., Mathieu, E., \& Ortiz-Ospina, E. (2024).
\emph{Our World in Data}. \url{https://ourworldindata.org/}

\bibitem[\citeproctext]{ref-tufte2001}
Tufte, E. R. (2001). \emph{The Visual Display of Quantitative
Information} (2. Aufl.). Graphics Press.

\bibitem[\citeproctext]{ref-wickham2019}
Wickham, H., Averick, M., Bryan, J., Chang, W., McGowan, L. D.,
François, R., Grolemund, G., Hayes, A., Henry, L., Hester, J., Kuhn, M.,
Pedersen, T. L., Miller, E., Bache, S. M., Müller, K., Ooms, J.,
Robinson, D., Seidel, D. P., Spinu, V., \ldots{} Yutani, H. (2019).
Welcome to the Tidyverse. \emph{Journal of Open Source Software},
\emph{4}(43), 1686. \url{https://doi.org/10.21105/joss.01686}

\end{CSLReferences}

\chapter{3a --- Inhaltsverzeichnis und
Literatur}\label{a-inhaltsverzeichnis-und-literatur}

\subsubsection{Was du nach diesem Kapitel
kannst}\label{was-du-nach-diesem-kapitel-kannst-2}

\begin{itemize}
\tightlist
\item
  ein \textbf{automatisches Inhaltsverzeichnis} in deinem Skript
  einschalten,
\item
  eine \textbf{Zotero-Bibliothek als BibTeX exportieren} und ins Projekt
  einbinden,
\item
  mit \texttt{{[}@key{]}} zitieren und ein \textbf{automatisches
  Literaturverzeichnis} am Kapitelende erzeugen.
\end{itemize}

\section{Worum es im Kern geht}\label{worum-es-im-kern-geht-1}

Ein gutes Skript orientiert die Lesenden zwei Mal: oben durch das
Inhaltsverzeichnis, unten durch die Literaturliste. Beides automatisch
zu erzeugen ist mehr als Komfort --- es vermeidet die typischen
Fehlerquellen (vergessene Quellen, falsche Seitenzahlen, doppelte
Einträge), die in Hand-pflege unweigerlich entstehen. Die Verbindung von
Zotero und Quarto regelt das mit zwei Konfigurationszeilen.

\section{Inhaltsverzeichnis --- drei
Zeilen}\label{inhaltsverzeichnis-drei-zeilen}

In \texttt{\_quarto.yml}:

\begin{Shaded}
\begin{Highlighting}[]
\FunctionTok{format}\KeywordTok{:}
\AttributeTok{  }\FunctionTok{html}\KeywordTok{:}
\AttributeTok{    }\FunctionTok{toc}\KeywordTok{:}\AttributeTok{ }\CharTok{true}
\AttributeTok{    }\FunctionTok{toc{-}depth}\KeywordTok{:}\AttributeTok{ }\DecValTok{3}
\AttributeTok{    }\FunctionTok{toc{-}title}\KeywordTok{:}\AttributeTok{ }\StringTok{"Inhalt"}
\end{Highlighting}
\end{Shaded}

Quarto baut das TOC automatisch aus den Überschriften. Tiefe drei heißt:
Kapitel, Abschnitt, Unterabschnitt --- alles, was tiefer geschachtelt
ist, taucht im TOC nicht auf.

\section{Zotero anbinden}\label{zotero-anbinden}

Zotero exportiert deine Bibliothek als \texttt{.bib}-Datei, Quarto liest
sie ein, du zitierst per Schlüssel. Das \textbf{Better-BibTeX}-Plugin
sorgt dafür, dass die Schlüssel stabil bleiben (z. B.
\texttt{mayer2021}) und der Export sich automatisch aktualisiert, wenn
du in Zotero etwas änderst.

\subsection{Mach mit --- in fünf
Schritten}\label{mach-mit-in-fuxfcnf-schritten}

\begin{enumerate}
\def\labelenumi{\arabic{enumi}.}
\tightlist
\item
  \textbf{Better BibTeX} in Zotero installieren:
  \url{https://retorque.re/zotero-better-bibtex/}.
\item
  Eine kleine \textbf{Zotero-Sammlung} anlegen mit zwei oder drei
  Quellen aus deinem Strang.
\item
  Sammlung \textbf{rechtsklicken → Export Collection → Better BibTeX →
  Keep updated}.
\item
  Datei speichern als \texttt{references.bib} im Projektordner.
\item
  In \texttt{\_quarto.yml}:
\end{enumerate}

\begin{Shaded}
\begin{Highlighting}[]
\FunctionTok{bibliography}\KeywordTok{:}\AttributeTok{ references.bib}
\FunctionTok{csl}\KeywordTok{:}\AttributeTok{ https://www.zotero.org/styles/apa}
\FunctionTok{link{-}citations}\KeywordTok{:}\AttributeTok{ }\CharTok{true}
\end{Highlighting}
\end{Shaded}

\subsection{Zitieren}\label{zitieren}

Im Text:

\begin{Shaded}
\begin{Highlighting}[]
\NormalTok{Generative Lernaktivitäten verbessern das Behalten }\CommentTok{[}\OtherTok{@fiorella2016}\CommentTok{]}\NormalTok{.}
\NormalTok{Die zugrundeliegende Theorie geht auf Mayer }\CommentTok{[}\OtherTok{{-}@mayer2021}\CommentTok{]}\NormalTok{ zurück.}
\end{Highlighting}
\end{Shaded}

Quarto rendert das in APA und hängt am Kapitelende automatisch eine
Literaturliste an, sofern du dort einen \texttt{\#\ Literatur}-Heading
mit \texttt{:::\ \{\#refs\}} setzt --- wie in diesem Skript.

\begin{tcolorbox}[enhanced jigsaw, opacitybacktitle=0.6, leftrule=.75mm, title=\textcolor{quarto-callout-tip-color}{\faLightbulb}\hspace{0.5em}{Mach mit --- strangbezogen}, colframe=quarto-callout-tip-color-frame, arc=.35mm, toprule=.15mm, colbacktitle=quarto-callout-tip-color!10!white, bottomrule=.15mm, colback=white, rightrule=.15mm, breakable, toptitle=1mm, bottomtitle=1mm, opacityback=0, titlerule=0mm, coltitle=black, left=2mm]

Wähle aus deinem Strang \textbf{drei Quellen}, lege sie in Zotero an,
exportiere und zitiere sie in deinem Kapitel. Wenn du keinen passenden
Startpunkt hast, hier je eine Empfehlung:

\begin{itemize}
\tightlist
\item
  \textbf{Strang A --- Wissenschaftliche Einstellung:} Bergstrom \& West
  (\citeproc{ref-bergstrom2024}{2024}), dazu zwei eigene Ergänzungen.
\item
  \textbf{Strang B --- Recherche und Quellenqualität:} Bergstrom \& West
  (\citeproc{ref-bergstrom2024}{2024}), dazu eine Quelle zu
  Open-Access-Datenbanken (z. B. \url{https://doaj.org}) und eine zu
  Pre-Print-Servern.
\item
  \textbf{Strang C --- Transparent schreiben:} Bergstrom \& West
  (\citeproc{ref-bergstrom2024}{2024}), dazu eine Quelle zur
  TOP-Guideline-Initiative und eine zu Reproducible Research.
\end{itemize}

\end{tcolorbox}

\begin{tcolorbox}[enhanced jigsaw, opacitybacktitle=0.6, leftrule=.75mm, title=\textcolor{quarto-callout-warning-color}{\faExclamationTriangle}\hspace{0.5em}{Was, wenn \ldots{}}, colframe=quarto-callout-warning-color-frame, arc=.35mm, toprule=.15mm, colbacktitle=quarto-callout-warning-color!10!white, bottomrule=.15mm, colback=white, rightrule=.15mm, breakable, toptitle=1mm, bottomtitle=1mm, opacityback=0, titlerule=0mm, coltitle=black, left=2mm]

\begin{itemize}
\tightlist
\item
  \textbf{„Citation key not found''?} → BibTeX-Schlüssel im
  \texttt{.bib} prüfen, oder Better BibTeX neu auto-exportieren.
\item
  \textbf{Liste erscheint nicht?} → Am Kapitelende
  \texttt{\#\ Literatur} und darunter \texttt{:::\ \{\#refs\}} mit
  \texttt{:::} einfügen. Quarto füllt diesen Block.
\item
  \textbf{Stil falsch?} → CSL-URL prüfen oder eine eigene CSL-Datei aus
  \url{https://www.zotero.org/styles} herunterladen und lokal einbinden
  (\texttt{csl:\ my-style.csl}).
\item
  \textbf{Mehrere Schlüssel zugleich?} →
  \texttt{{[}@fiorella2016;\ @mayer2021{]}} mit Semikolon trennen.
\end{itemize}

\end{tcolorbox}

\section{Tutor}\label{tutor-2}

\begin{tcolorbox}[enhanced jigsaw, opacitybacktitle=0.6, leftrule=.75mm, title=\textcolor{quarto-callout-tip-color}{\faLightbulb}\hspace{0.5em}{Tutor öffnen}, colframe=quarto-callout-tip-color-frame, arc=.35mm, toprule=.15mm, colbacktitle=quarto-callout-tip-color!10!white, bottomrule=.15mm, colback=white, rightrule=.15mm, breakable, toptitle=1mm, bottomtitle=1mm, opacityback=0, titlerule=0mm, coltitle=black, left=2mm]

\href{https://chatgpt.com/g/g-69fec93348108191af1ffb8ce0cf6712-tutor-interaktive-skripte}{Tutor
„Interaktive Skripte'' →}

Vorschlag-Prompt für dieses Kapitel:

\begin{quote}
Ich möchte aus Zotero exportierte Quellen in einem Quarto-Skript
zitieren. Mein Stand: {[}Better BibTeX installiert / nicht
installiert{]}; {[}references.bib liegt im Projekt / fehlt noch{]}.
Erkläre mir den nächsten Schritt mit Lay-Erklärung vor dem Code. Frag,
wo ich gerade stehe.
\end{quote}

\end{tcolorbox}

\chapter*{Literatur}\label{literatur-2}
\addcontentsline{toc}{chapter}{Literatur}

\markboth{Literatur}{Literatur}

\phantomsection\label{refs}
\begin{CSLReferences}{1}{0}
\bibitem[\citeproctext]{ref-allaire2024}
Allaire, J. J., Teague, C., Scheidegger, C., Xie, Y., \& Dervieux, C.
(2024). \emph{Quarto: An Open-Source Scientific and Technical Publishing
System}. \url{https://quarto.org/}

\bibitem[\citeproctext]{ref-bergstrom2024}
Bergstrom, C. T., \& West, J. D. (2024). \emph{Calling Bullshit:
Wissenschaftliches Arbeiten und kritisches Denken}.
\url{https://publish.obsidian.md/wagp-calling-bullshit/}

\bibitem[\citeproctext]{ref-bryan2018}
Bryan, J. (2018). Excuse Me, Do You Have a Moment to Talk About Version
Control? \emph{The American Statistician}, \emph{72}(1), 20--27.
\url{https://doi.org/10.1080/00031305.2017.1399928}

\bibitem[\citeproctext]{ref-fiorella2016}
Fiorella, L., \& Mayer, R. E. (2016). Eight Ways to Promote Generative
Learning. \emph{Educational Psychology Review}, \emph{28}(4), 717--741.
\url{https://doi.org/10.1007/s10648-015-9348-9}

\bibitem[\citeproctext]{ref-healy2018}
Healy, K. (2018). \emph{Data Visualization: A Practical Introduction}.
Princeton University Press.

\bibitem[\citeproctext]{ref-karpicke2011}
Karpicke, J. D., \& Blunt, J. R. (2011). Retrieval Practice Produces
More Learning than Elaborative Studying with Concept Mapping.
\emph{Science}, \emph{331}(6018), 772--775.
\url{https://doi.org/10.1126/science.1199327}

\bibitem[\citeproctext]{ref-kasneci2023}
Kasneci, E., Sessler, K., Küchemann, S., Bannert, M., Dementieva, D.,
Fischer, F., Gasser, U., Groh, G., Günnemann, S., Hüllermeier, E.,
Krusche, S., Kutyniok, G., Michaeli, T., Nerdel, C., Pfeffer, J.,
Poquet, O., Sailer, M., Schmidt, A., Seidel, T., \ldots{} Kasneci, G.
(2023). {ChatGPT} for Good? On Opportunities and Challenges of Large
Language Models for Education. \emph{Learning and Individual
Differences}, \emph{103}, 102274.
\url{https://doi.org/10.1016/j.lindif.2023.102274}

\bibitem[\citeproctext]{ref-lakens2022}
Lakens, D. (2022). \emph{Improving Your Statistical Inferences}.
\url{https://doi.org/10.5281/zenodo.6409077}

\bibitem[\citeproctext]{ref-mayer2021}
Mayer, R. E. (2021). \emph{Multimedia Learning} (3. Aufl.). Cambridge
University Press. \url{https://doi.org/10.1017/9781316941355}

\bibitem[\citeproctext]{ref-mollick2023}
Mollick, E. R., \& Mollick, L. (2023). Assigning {AI}: Seven Approaches
for Students, with Prompts. \emph{SSRN Working Paper}.
\url{https://doi.org/10.2139/ssrn.4475995}

\bibitem[\citeproctext]{ref-roediger2006}
Roediger, H. L., \& Karpicke, J. D. (2006). Test-Enhanced Learning:
Taking Memory Tests Improves Long-Term Retention. \emph{Psychological
Science}, \emph{17}(3), 249--255.
\url{https://doi.org/10.1111/j.1467-9280.2006.01693.x}

\bibitem[\citeproctext]{ref-owid2024}
Roser, M., Ritchie, H., Mathieu, E., \& Ortiz-Ospina, E. (2024).
\emph{Our World in Data}. \url{https://ourworldindata.org/}

\bibitem[\citeproctext]{ref-tufte2001}
Tufte, E. R. (2001). \emph{The Visual Display of Quantitative
Information} (2. Aufl.). Graphics Press.

\bibitem[\citeproctext]{ref-wickham2019}
Wickham, H., Averick, M., Bryan, J., Chang, W., McGowan, L. D.,
François, R., Grolemund, G., Hayes, A., Henry, L., Hester, J., Kuhn, M.,
Pedersen, T. L., Miller, E., Bache, S. M., Müller, K., Ooms, J.,
Robinson, D., Seidel, D. P., Spinu, V., \ldots{} Yutani, H. (2019).
Welcome to the Tidyverse. \emph{Journal of Open Source Software},
\emph{4}(43), 1686. \url{https://doi.org/10.21105/joss.01686}

\end{CSLReferences}

\chapter{3b --- TH-Köln-Stil einbauen}\label{b-th-kuxf6ln-stil-einbauen}

\subsubsection{Was du nach diesem Kapitel
kannst}\label{was-du-nach-diesem-kapitel-kannst-3}

\begin{itemize}
\tightlist
\item
  die Hochschulfarben über \texttt{\_brand.yml} setzen,
\item
  per CSS scharfkantige Callouts, Tabellen und Code-Blöcke gestalten,
\item
  das TH-Köln-Logo einbinden und prüfen, dass es überall korrekt
  erscheint.
\end{itemize}

\section{Worum es im Kern geht}\label{worum-es-im-kern-geht-2}

Gestaltung ist nicht Dekoration, sondern Lese-Hilfe. Ein konsistentes
Layout reduziert kognitive Last (\citeproc{ref-mayer2021}{Mayer, 2021})
--- die Lernenden müssen nicht für jedes neue Element wieder verstehen,
was wo wichtig ist. Tufte (\citeproc{ref-tufte2001}{2001}) formuliert
dasselbe Prinzip für Datenvisualisierung: maximaler Informationsgehalt
bei minimalem Tinte-zu-Daten-Verhältnis. Beides spricht für
Zurückhaltung im Stil --- wenige Farben, klare Hierarchien, kein
Schmuck.

\section{Wie Quarto Aussehen
handhabt}\label{wie-quarto-aussehen-handhabt}

Drei Stellschrauben, in dieser Reihenfolge:

\begin{enumerate}
\def\labelenumi{\arabic{enumi}.}
\tightlist
\item
  \textbf{\texttt{\_brand.yml}} --- die Hochschulfarben und -schriften
  an einem Ort.
\item
  \textbf{\texttt{styles.scss}} --- alles, was \texttt{\_brand.yml}
  nicht abdeckt (Layout, Akzente, eigene Boxen).
\item
  \textbf{\texttt{format.html.theme}} in \texttt{\_quarto.yml} --- wählt
  das Bootstrap-Basistheme.
\end{enumerate}

\section{Die TH-Köln-Palette in
Stichworten}\label{die-th-kuxf6ln-palette-in-stichworten}

Drei Farben, sonst nur Grautöne. Alles andere ist Verstoß gegen das
Corporate Design.

\begin{longtable}[]{@{}lll@{}}
\toprule\noalign{}
Rolle & Hex & Wo? \\
\midrule\noalign{}
\endhead
\bottomrule\noalign{}
\endlastfoot
TH Red & \texttt{\#c81e0f} & Links, Akzente, wichtige Hervorhebungen \\
TH Magenta & \texttt{\#b43092} & Strukturlinien, Callout-Ränder \\
TH Orange & \texttt{\#ea5a00} & Sehr sparsam --- Warnungen, Verlauf \\
Ink & \texttt{\#111111} & Fließtext (weicheres Schwarz) \\
Paper & \texttt{\#ffffff} & Hintergründe, immer \\
\end{longtable}

\section{\texorpdfstring{Mach mit --- \texttt{\_brand.yml}
kopieren}{Mach mit --- \_brand.yml kopieren}}\label{mach-mit-_brand.yml-kopieren}

Im Repository des Workshops liegt eine fertige \texttt{\_brand.yml} mit
der TH-Köln-Palette. Kopiere sie in dein Projektverzeichnis. Beim
nächsten \texttt{quarto\ render} sind Links rot, Überschriften schwarz,
Akzente magenta.

\begin{Shaded}
\begin{Highlighting}[]
\FunctionTok{color}\KeywordTok{:}
\AttributeTok{  }\FunctionTok{palette}\KeywordTok{:}
\AttributeTok{    }\FunctionTok{th{-}red}\KeywordTok{:}\AttributeTok{ }\StringTok{"\#c81e0f"}
\AttributeTok{    }\FunctionTok{th{-}magenta}\KeywordTok{:}\AttributeTok{ }\StringTok{"\#b43092"}
\AttributeTok{    }\FunctionTok{th{-}orange}\KeywordTok{:}\AttributeTok{ }\StringTok{"\#ea5a00"}
\AttributeTok{  }\FunctionTok{primary}\KeywordTok{:}\AttributeTok{ th{-}red}
\AttributeTok{  }\FunctionTok{secondary}\KeywordTok{:}\AttributeTok{ th{-}magenta}
\end{Highlighting}
\end{Shaded}

\section{SCSS-Ergänzungen}\label{scss-erguxe4nzungen}

Drei Dinge, die \texttt{\_brand.yml} nicht regelt: scharfe Ecken
(„kantig''), die magenta H2-Unterstreichung und die TH-Farbleiste am
Seitenanfang. Quarto verarbeitet eigene Stile als \textbf{SCSS} (Sass)
--- eine Erweiterung von CSS, die mit Bootstrap zusammenarbeitet. Du
brauchst dafür kein Sass-Wissen: jede gültige CSS-Regel ist auch
gültiges SCSS.

\begin{tcolorbox}[enhanced jigsaw, opacitybacktitle=0.6, leftrule=.75mm, title=\textcolor{quarto-callout-warning-color}{\faExclamationTriangle}\hspace{0.5em}{Warnung}, colframe=quarto-callout-warning-color-frame, arc=.35mm, toprule=.15mm, colbacktitle=quarto-callout-warning-color!10!white, bottomrule=.15mm, colback=white, rightrule=.15mm, breakable, toptitle=1mm, bottomtitle=1mm, opacityback=0, titlerule=0mm, coltitle=black, left=2mm]

Die Datei muss \texttt{styles.scss} heißen (nicht \texttt{.css}) und am
Anfang einen \textbf{Layer-Marker} enthalten --- sonst bricht
\texttt{quarto\ render} mit einer Fehlermeldung „doesn't contain at
least one layer boundary'' ab.

\end{tcolorbox}

Lege \texttt{styles.scss} an und binde sie in \texttt{\_quarto.yml} ein:

\begin{Shaded}
\begin{Highlighting}[]
\FunctionTok{format}\KeywordTok{:}
\AttributeTok{  }\FunctionTok{html}\KeywordTok{:}
\AttributeTok{    }\FunctionTok{theme}\KeywordTok{:}
\AttributeTok{      }\KeywordTok{{-}}\AttributeTok{ cosmo}
\AttributeTok{      }\KeywordTok{{-}}\AttributeTok{ styles.scss}
\end{Highlighting}
\end{Shaded}

In \texttt{styles.scss} (Auszug --- Layer-Marker am Anfang ist Pflicht):

\begin{Shaded}
\begin{Highlighting}[]
\CommentTok{/*{-}{-} scss:rules {-}{-}*/}

\InformationTok{:root}\NormalTok{ \{}
  \VariableTok{{-}{-}th{-}red}\NormalTok{: }\ConstantTok{\#c81e0f}\OperatorTok{;}
  \VariableTok{{-}{-}th{-}magenta}\NormalTok{: }\ConstantTok{\#b43092}\OperatorTok{;}
\NormalTok{\}}

\NormalTok{h1}\OperatorTok{,}\NormalTok{ h2}\OperatorTok{,}\NormalTok{ h3 \{ }\KeywordTok{border{-}radius}\NormalTok{: }\DecValTok{0} \AttributeTok{!important}\OperatorTok{;} \KeywordTok{font{-}weight}\NormalTok{: }\DecValTok{600}\OperatorTok{;}\NormalTok{ \}}
\NormalTok{h2 \{ }\KeywordTok{border{-}bottom}\NormalTok{: }\DecValTok{2}\DataTypeTok{px} \DecValTok{solid} \FunctionTok{rgba(}\DecValTok{180}\OperatorTok{,} \DecValTok{48}\OperatorTok{,} \DecValTok{146}\OperatorTok{,} \DecValTok{0.35}\FunctionTok{)}\OperatorTok{;} \KeywordTok{padding{-}bottom}\NormalTok{: }\DecValTok{0.3}\DataTypeTok{rem}\OperatorTok{;}\NormalTok{ \}}
\FunctionTok{.callout}\NormalTok{ \{ }\KeywordTok{border{-}radius}\NormalTok{: }\DecValTok{0} \AttributeTok{!important}\OperatorTok{;}\NormalTok{ \}}
\NormalTok{pre}\OperatorTok{,}\NormalTok{ code \{ }\KeywordTok{border{-}radius}\NormalTok{: }\DecValTok{0} \AttributeTok{!important}\OperatorTok{;}\NormalTok{ \}}
\end{Highlighting}
\end{Shaded}

\begin{tcolorbox}[enhanced jigsaw, opacitybacktitle=0.6, leftrule=.75mm, title=\textcolor{quarto-callout-tip-color}{\faLightbulb}\hspace{0.5em}{Mach mit --- Vorher/Nachher}, colframe=quarto-callout-tip-color-frame, arc=.35mm, toprule=.15mm, colbacktitle=quarto-callout-tip-color!10!white, bottomrule=.15mm, colback=white, rightrule=.15mm, breakable, toptitle=1mm, bottomtitle=1mm, opacityback=0, titlerule=0mm, coltitle=black, left=2mm]

Render dein Skript einmal \textbf{vor} und einmal \textbf{nach} dem
Einbau von \texttt{\_brand.yml}. Vergleich beide Versionen
nebeneinander. Was ändert sich? Was nicht? Schreib dir drei
Beobachtungen auf --- wir nutzen sie im Showcase am Ende.

\end{tcolorbox}

\section{Logo einbinden}\label{logo-einbinden}

\begin{enumerate}
\def\labelenumi{\arabic{enumi}.}
\tightlist
\item
  Lade das TH-Köln-Logo (PNG, transparent) in
  \texttt{images/th-koeln-logo.png}.
\item
  In \texttt{\_brand.yml}:
\end{enumerate}

\begin{Shaded}
\begin{Highlighting}[]
\FunctionTok{logo}\KeywordTok{:}
\AttributeTok{  }\FunctionTok{small}\KeywordTok{:}\AttributeTok{ images/th{-}koeln{-}logo.png}
\AttributeTok{  }\FunctionTok{medium}\KeywordTok{:}\AttributeTok{ images/th{-}koeln{-}logo.png}
\end{Highlighting}
\end{Shaded}

\begin{enumerate}
\def\labelenumi{\arabic{enumi}.}
\setcounter{enumi}{2}
\tightlist
\item
  Render und prüfen.
\end{enumerate}

\begin{tcolorbox}[enhanced jigsaw, opacitybacktitle=0.6, leftrule=.75mm, title=\textcolor{quarto-callout-warning-color}{\faExclamationTriangle}\hspace{0.5em}{Was, wenn \ldots{}}, colframe=quarto-callout-warning-color-frame, arc=.35mm, toprule=.15mm, colbacktitle=quarto-callout-warning-color!10!white, bottomrule=.15mm, colback=white, rightrule=.15mm, breakable, toptitle=1mm, bottomtitle=1mm, opacityback=0, titlerule=0mm, coltitle=black, left=2mm]

\begin{itemize}
\tightlist
\item
  \textbf{Farben tauchen nicht auf?} → Cache leeren
  (\texttt{quarto\ render\ -\/-cache-refresh}) oder den Ordner
  \texttt{\_freeze/} löschen.
\item
  \textbf{CSS wird nicht geladen?} → Pfad in \texttt{\_quarto.yml}
  prüfen, relative Pfade ohne führenden Slash.
\item
  \textbf{Logo wird nicht angezeigt?} → PNG-Datei muss tatsächlich im
  \texttt{images/}-Ordner liegen, nicht nur in \texttt{\_brand.yml}
  referenziert sein.
\item
  \textbf{Farben sehen ausgewaschen aus?} → Eventuell ein
  Bootstrap-Theme aktiv, das die Variablen überschreibt.
  \texttt{theme:\ {[}cosmo,\ styles.scss{]}} in dieser Reihenfolge
  halten --- die SCSS-Datei muss zuletzt geladen werden.
\item
  \textbf{„doesn't contain at least one layer boundary''} beim Render? →
  Am Anfang von \texttt{styles.scss} muss
  \texttt{/*-\/-\ scss:rules\ -\/-*/} stehen.
\end{itemize}

\end{tcolorbox}

\section{Tutor}\label{tutor-3}

\begin{tcolorbox}[enhanced jigsaw, opacitybacktitle=0.6, leftrule=.75mm, title=\textcolor{quarto-callout-tip-color}{\faLightbulb}\hspace{0.5em}{Tutor öffnen}, colframe=quarto-callout-tip-color-frame, arc=.35mm, toprule=.15mm, colbacktitle=quarto-callout-tip-color!10!white, bottomrule=.15mm, colback=white, rightrule=.15mm, breakable, toptitle=1mm, bottomtitle=1mm, opacityback=0, titlerule=0mm, coltitle=black, left=2mm]

\href{https://chatgpt.com/g/g-69fec93348108191af1ffb8ce0cf6712-tutor-interaktive-skripte}{Tutor
„Interaktive Skripte'' →}

Vorschlag-Prompt für dieses Kapitel:

\begin{quote}
Ich möchte mein Quarto-Skript im TH-Köln-Stil gestalten (rot \#c81e0f,
magenta \#b43092, orange \#ea5a00, kantige Ecken, Arial). Mein Stand:
{[}\_brand.yml vorhanden / fehlt{]}; {[}styles.scss vorhanden /
fehlt{]}. Was ist mein nächster Schritt? Lay-Erklärung zuerst, dann ein
konkreter Code-Schnipsel.
\end{quote}

\end{tcolorbox}

\chapter*{Literatur}\label{literatur-3}
\addcontentsline{toc}{chapter}{Literatur}

\markboth{Literatur}{Literatur}

\phantomsection\label{refs}
\begin{CSLReferences}{1}{0}
\bibitem[\citeproctext]{ref-allaire2024}
Allaire, J. J., Teague, C., Scheidegger, C., Xie, Y., \& Dervieux, C.
(2024). \emph{Quarto: An Open-Source Scientific and Technical Publishing
System}. \url{https://quarto.org/}

\bibitem[\citeproctext]{ref-bergstrom2024}
Bergstrom, C. T., \& West, J. D. (2024). \emph{Calling Bullshit:
Wissenschaftliches Arbeiten und kritisches Denken}.
\url{https://publish.obsidian.md/wagp-calling-bullshit/}

\bibitem[\citeproctext]{ref-bryan2018}
Bryan, J. (2018). Excuse Me, Do You Have a Moment to Talk About Version
Control? \emph{The American Statistician}, \emph{72}(1), 20--27.
\url{https://doi.org/10.1080/00031305.2017.1399928}

\bibitem[\citeproctext]{ref-fiorella2016}
Fiorella, L., \& Mayer, R. E. (2016). Eight Ways to Promote Generative
Learning. \emph{Educational Psychology Review}, \emph{28}(4), 717--741.
\url{https://doi.org/10.1007/s10648-015-9348-9}

\bibitem[\citeproctext]{ref-healy2018}
Healy, K. (2018). \emph{Data Visualization: A Practical Introduction}.
Princeton University Press.

\bibitem[\citeproctext]{ref-karpicke2011}
Karpicke, J. D., \& Blunt, J. R. (2011). Retrieval Practice Produces
More Learning than Elaborative Studying with Concept Mapping.
\emph{Science}, \emph{331}(6018), 772--775.
\url{https://doi.org/10.1126/science.1199327}

\bibitem[\citeproctext]{ref-kasneci2023}
Kasneci, E., Sessler, K., Küchemann, S., Bannert, M., Dementieva, D.,
Fischer, F., Gasser, U., Groh, G., Günnemann, S., Hüllermeier, E.,
Krusche, S., Kutyniok, G., Michaeli, T., Nerdel, C., Pfeffer, J.,
Poquet, O., Sailer, M., Schmidt, A., Seidel, T., \ldots{} Kasneci, G.
(2023). {ChatGPT} for Good? On Opportunities and Challenges of Large
Language Models for Education. \emph{Learning and Individual
Differences}, \emph{103}, 102274.
\url{https://doi.org/10.1016/j.lindif.2023.102274}

\bibitem[\citeproctext]{ref-lakens2022}
Lakens, D. (2022). \emph{Improving Your Statistical Inferences}.
\url{https://doi.org/10.5281/zenodo.6409077}

\bibitem[\citeproctext]{ref-mayer2021}
Mayer, R. E. (2021). \emph{Multimedia Learning} (3. Aufl.). Cambridge
University Press. \url{https://doi.org/10.1017/9781316941355}

\bibitem[\citeproctext]{ref-mollick2023}
Mollick, E. R., \& Mollick, L. (2023). Assigning {AI}: Seven Approaches
for Students, with Prompts. \emph{SSRN Working Paper}.
\url{https://doi.org/10.2139/ssrn.4475995}

\bibitem[\citeproctext]{ref-roediger2006}
Roediger, H. L., \& Karpicke, J. D. (2006). Test-Enhanced Learning:
Taking Memory Tests Improves Long-Term Retention. \emph{Psychological
Science}, \emph{17}(3), 249--255.
\url{https://doi.org/10.1111/j.1467-9280.2006.01693.x}

\bibitem[\citeproctext]{ref-owid2024}
Roser, M., Ritchie, H., Mathieu, E., \& Ortiz-Ospina, E. (2024).
\emph{Our World in Data}. \url{https://ourworldindata.org/}

\bibitem[\citeproctext]{ref-tufte2001}
Tufte, E. R. (2001). \emph{The Visual Display of Quantitative
Information} (2. Aufl.). Graphics Press.

\bibitem[\citeproctext]{ref-wickham2019}
Wickham, H., Averick, M., Bryan, J., Chang, W., McGowan, L. D.,
François, R., Grolemund, G., Hayes, A., Henry, L., Hester, J., Kuhn, M.,
Pedersen, T. L., Miller, E., Bache, S. M., Müller, K., Ooms, J.,
Robinson, D., Seidel, D. P., Spinu, V., \ldots{} Yutani, H. (2019).
Welcome to the Tidyverse. \emph{Journal of Open Source Software},
\emph{4}(43), 1686. \url{https://doi.org/10.21105/joss.01686}

\end{CSLReferences}

\chapter{4 --- Interaktive Diagramme}\label{interaktive-diagramme}

\subsubsection{Was du nach diesem Kapitel
kannst}\label{was-du-nach-diesem-kapitel-kannst-4}

\begin{itemize}
\tightlist
\item
  ein \textbf{interaktives Diagramm} mit echten Daten in dein Skript
  einbauen,
\item
  Daten direkt aus dem Netz laden (z. B. \textbf{Our World in Data}),
\item
  erklären, \textbf{wann} sich Interaktivität didaktisch lohnt und wann
  nicht.
\end{itemize}

\section{Worum es im Kern geht}\label{worum-es-im-kern-geht-3}

Diagramme tun in einem Skript zwei Dinge auf einmal --- sie zeigen Daten
und sie laden zur Interpretation ein. Tufte
(\citeproc{ref-tufte2001}{2001}) formuliert die Maxime: maximale
Information bei minimalem Schmuck. Healy
(\citeproc{ref-healy2018}{2018}) ergänzt für die digitale Lehre:
Interaktivität rechtfertigt sich nur, wenn die Lernenden tatsächlich
etwas mit ihr anfangen --- verschieben, ein- und ausblenden,
vergleichen. Andernfalls reicht ein PNG.

\section{Drei Wege, ein Diagramm
einzubauen}\label{drei-wege-ein-diagramm-einzubauen}

\begin{longtable}[]{@{}
  >{\raggedright\arraybackslash}p{(\linewidth - 4\tabcolsep) * \real{0.3333}}
  >{\raggedright\arraybackslash}p{(\linewidth - 4\tabcolsep) * \real{0.3333}}
  >{\raggedright\arraybackslash}p{(\linewidth - 4\tabcolsep) * \real{0.3333}}@{}}
\toprule\noalign{}
\begin{minipage}[b]{\linewidth}\raggedright
Weg
\end{minipage} & \begin{minipage}[b]{\linewidth}\raggedright
Werkzeug
\end{minipage} & \begin{minipage}[b]{\linewidth}\raggedright
Wann passend?
\end{minipage} \\
\midrule\noalign{}
\endhead
\bottomrule\noalign{}
\endlastfoot
\textbf{R + ggplot2 + plotly} & R-Code-Chunk & Du kennst R oder willst
es lernen. Volle Kontrolle, schöne Grafiken
(\citeproc{ref-wickham2019}{Wickham et al., 2019}). \\
\textbf{Observable JS} & OJS-Chunk in Quarto & Du willst dich auf
JavaScript einlassen. Sehr schnell, modern. \\
\textbf{iFrame mit Datawrapper / Flourish} & externer Editor, Embed-Code
& Du willst gar nicht programmieren. \\
\end{longtable}

Wir gehen heute den ersten Weg --- er ist mit RStudio am robustesten und
nutzt die R-Pakete, die du wahrscheinlich ohnehin im Werkzeugkasten
haben willst.

\section{Mach mit --- ein interaktives Diagramm in 10
Minuten}\label{mach-mit-ein-interaktives-diagramm-in-10-minuten}

Wir laden CO₂-Daten von Our World in Data (\citeproc{ref-owid2024}{Roser
et al., 2024}) und zeichnen sie als interaktive Linie. Drei
Strang-Varianten --- jede passt auf das Stoffgebiet, das du gewählt
hast.

\subsubsection{Strang A --- Wissenschaftliche Einstellung}

Frage: \emph{Wie haben sich Schätzungen zur globalen Lebenserwartung
über die Zeit verändert?} Diese Frage zeigt, dass „Fakten'' historisch
sind --- sie verändern sich mit besserer Datenlage.

\begin{Shaded}
\begin{Highlighting}[]
\InformationTok{\textasciigrave{}\textasciigrave{}\textasciigrave{}\{r\}}
\CommentTok{\#| label: fig{-}le}
\CommentTok{\#| fig{-}cap: "Lebenserwartung bei Geburt, ausgewählte Länder. Daten: Our World in Data (Roser et al., 2024)."}
\CommentTok{\#| message: false}

\FunctionTok{library}\NormalTok{(plotly)}
\FunctionTok{library}\NormalTok{(dplyr)}
\FunctionTok{library}\NormalTok{(readr)}

\NormalTok{owid }\OtherTok{\textless{}{-}} \FunctionTok{read\_csv}\NormalTok{(}\StringTok{"https://ourworldindata.org/grapher/life{-}expectancy.csv?v=1"}\NormalTok{)}

\NormalTok{ausgewaehlt }\OtherTok{\textless{}{-}}\NormalTok{ owid }\SpecialCharTok{|\textgreater{}}
  \FunctionTok{filter}\NormalTok{(Entity }\SpecialCharTok{\%in\%} \FunctionTok{c}\NormalTok{(}\StringTok{"Germany"}\NormalTok{, }\StringTok{"Japan"}\NormalTok{, }\StringTok{"World"}\NormalTok{),}
\NormalTok{         Year }\SpecialCharTok{\textgreater{}=} \DecValTok{1950}\NormalTok{)}

\FunctionTok{plot\_ly}\NormalTok{(ausgewaehlt,}
        \AttributeTok{x =} \SpecialCharTok{\textasciitilde{}}\NormalTok{Year, }\AttributeTok{y =} \SpecialCharTok{\textasciitilde{}}\StringTok{\textasciigrave{}}\AttributeTok{Life expectancy at birth (historical)}\StringTok{\textasciigrave{}}\NormalTok{,}
        \AttributeTok{color =} \SpecialCharTok{\textasciitilde{}}\NormalTok{Entity, }\AttributeTok{type =} \StringTok{"scatter"}\NormalTok{, }\AttributeTok{mode =} \StringTok{"lines"}\NormalTok{,}
        \AttributeTok{colors =} \FunctionTok{c}\NormalTok{(}\StringTok{"\#c81e0f"}\NormalTok{, }\StringTok{"\#b43092"}\NormalTok{, }\StringTok{"\#ea5a00"}\NormalTok{)) }\SpecialCharTok{|\textgreater{}}
  \FunctionTok{layout}\NormalTok{(}\AttributeTok{yaxis =} \FunctionTok{list}\NormalTok{(}\AttributeTok{title =} \StringTok{"Lebenserwartung (Jahre)"}\NormalTok{),}
         \AttributeTok{xaxis =} \FunctionTok{list}\NormalTok{(}\AttributeTok{title =} \StringTok{""}\NormalTok{))}
\InformationTok{\textasciigrave{}\textasciigrave{}\textasciigrave{}}
\end{Highlighting}
\end{Shaded}

\subsubsection{Strang B --- Recherche und Quellenqualität}

Frage: \emph{Wie viele wissenschaftliche Publikationen erscheinen pro
Jahr?} Antwort führt zur Notwendigkeit, Quellen zu filtern.

\begin{Shaded}
\begin{Highlighting}[]
\InformationTok{\textasciigrave{}\textasciigrave{}\textasciigrave{}\{r\}}
\CommentTok{\#| label: fig{-}pub}
\CommentTok{\#| fig{-}cap: "Wissenschaftliche Publikationen weltweit. Daten: Our World in Data."}
\CommentTok{\#| message: false}

\FunctionTok{library}\NormalTok{(plotly)}
\FunctionTok{library}\NormalTok{(dplyr)}
\FunctionTok{library}\NormalTok{(readr)}

\NormalTok{owid }\OtherTok{\textless{}{-}} \FunctionTok{read\_csv}\NormalTok{(}\StringTok{"https://ourworldindata.org/grapher/scientific{-}publications.csv?v=1"}\NormalTok{)}

\FunctionTok{plot\_ly}\NormalTok{(owid }\SpecialCharTok{|\textgreater{}} \FunctionTok{filter}\NormalTok{(Year }\SpecialCharTok{\textgreater{}=} \DecValTok{1990}\NormalTok{),}
        \AttributeTok{x =} \SpecialCharTok{\textasciitilde{}}\NormalTok{Year, }\AttributeTok{y =} \SpecialCharTok{\textasciitilde{}}\StringTok{\textasciigrave{}}\AttributeTok{Articles}\StringTok{\textasciigrave{}}\NormalTok{,}
        \AttributeTok{color =} \SpecialCharTok{\textasciitilde{}}\NormalTok{Entity, }\AttributeTok{type =} \StringTok{"scatter"}\NormalTok{, }\AttributeTok{mode =} \StringTok{"lines"}\NormalTok{,}
        \AttributeTok{colors =} \FunctionTok{c}\NormalTok{(}\StringTok{"\#c81e0f"}\NormalTok{, }\StringTok{"\#b43092"}\NormalTok{, }\StringTok{"\#ea5a00"}\NormalTok{)) }\SpecialCharTok{|\textgreater{}}
  \FunctionTok{layout}\NormalTok{(}\AttributeTok{yaxis =} \FunctionTok{list}\NormalTok{(}\AttributeTok{title =} \StringTok{"Publikationen pro Jahr"}\NormalTok{),}
         \AttributeTok{xaxis =} \FunctionTok{list}\NormalTok{(}\AttributeTok{title =} \StringTok{""}\NormalTok{))}
\InformationTok{\textasciigrave{}\textasciigrave{}\textasciigrave{}}
\end{Highlighting}
\end{Shaded}

\subsubsection{Strang C --- Transparent schreiben}

Frage: \emph{Wie hat sich der Anteil registrierter Studien
(Pre-Registration) entwickelt?} Datenquelle hier ggf. selbst
recherchieren --- Beispiel-Code unten lädt einen kleinen lokalen CSV.

\begin{Shaded}
\begin{Highlighting}[]
\InformationTok{\textasciigrave{}\textasciigrave{}\textasciigrave{}\{r\}}
\CommentTok{\#| label: fig{-}prereg}
\CommentTok{\#| fig{-}cap: "Anzahl pre{-}registrierter Studien auf OSF. Daten: eigener Auszug."}
\CommentTok{\#| message: false}

\FunctionTok{library}\NormalTok{(plotly)}
\FunctionTok{library}\NormalTok{(readr)}

\NormalTok{prereg }\OtherTok{\textless{}{-}} \FunctionTok{read\_csv}\NormalTok{(}\StringTok{"data/osf{-}preregistrations.csv"}\NormalTok{)}

\FunctionTok{plot\_ly}\NormalTok{(prereg, }\AttributeTok{x =} \SpecialCharTok{\textasciitilde{}}\NormalTok{Jahr, }\AttributeTok{y =} \SpecialCharTok{\textasciitilde{}}\NormalTok{Anzahl, }\AttributeTok{type =} \StringTok{"bar"}\NormalTok{,}
        \AttributeTok{marker =} \FunctionTok{list}\NormalTok{(}\AttributeTok{color =} \StringTok{"\#c81e0f"}\NormalTok{)) }\SpecialCharTok{|\textgreater{}}
  \FunctionTok{layout}\NormalTok{(}\AttributeTok{yaxis =} \FunctionTok{list}\NormalTok{(}\AttributeTok{title =} \StringTok{"Pre{-}Registrierungen"}\NormalTok{),}
         \AttributeTok{xaxis =} \FunctionTok{list}\NormalTok{(}\AttributeTok{title =} \StringTok{""}\NormalTok{))}
\InformationTok{\textasciigrave{}\textasciigrave{}\textasciigrave{}}
\end{Highlighting}
\end{Shaded}

Im Browser kannst du jetzt mit der Maus über die Linien fahren, Länder
ein- und ausschalten, in Bereiche zoomen. Das Bild im Skript ist live.

\section{\texorpdfstring{Wann \textbf{kein} interaktives
Diagramm?}{Wann kein interaktives Diagramm?}}\label{wann-kein-interaktives-diagramm}

Drei Situationen, in denen ein statisches PNG genauso gut wirkt:

\begin{itemize}
\tightlist
\item
  \textbf{Reine Illustration} --- wenn die Lernenden nichts ändern
  sollen, sondern nur etwas „sehen'', reicht ein Bild
  (\citeproc{ref-tufte2001}{Tufte, 2001}).
\item
  \textbf{Print-Version} wichtig --- interaktive Plots werden im PDF zu
  statischen Screenshots; wenn das PDF die Hauptform ist, plane direkt
  für Print.
\item
  \textbf{Performance} --- sehr große Datensätze (\textgreater100.000
  Punkte) machen plotly träge. Dann besser aggregieren oder ein
  statisches ggplot exportieren.
\end{itemize}

\begin{tcolorbox}[enhanced jigsaw, opacitybacktitle=0.6, leftrule=.75mm, title=\textcolor{quarto-callout-tip-color}{\faLightbulb}\hspace{0.5em}{Mach mit}, colframe=quarto-callout-tip-color-frame, arc=.35mm, toprule=.15mm, colbacktitle=quarto-callout-tip-color!10!white, bottomrule=.15mm, colback=white, rightrule=.15mm, breakable, toptitle=1mm, bottomtitle=1mm, opacityback=0, titlerule=0mm, coltitle=black, left=2mm]

Tausche im Code zwei Länder gegen welche aus, die für deinen Strang
interessant sind. Render, schau das Ergebnis im Browser an, push.

\end{tcolorbox}

\begin{tcolorbox}[enhanced jigsaw, opacitybacktitle=0.6, leftrule=.75mm, title=\textcolor{quarto-callout-warning-color}{\faExclamationTriangle}\hspace{0.5em}{Was, wenn \ldots{}}, colframe=quarto-callout-warning-color-frame, arc=.35mm, toprule=.15mm, colbacktitle=quarto-callout-warning-color!10!white, bottomrule=.15mm, colback=white, rightrule=.15mm, breakable, toptitle=1mm, bottomtitle=1mm, opacityback=0, titlerule=0mm, coltitle=black, left=2mm]

\begin{itemize}
\tightlist
\item
  \textbf{\texttt{could\ not\ find\ function\ "plot\_ly"}} →
  \texttt{install.packages("plotly")} in der R-Konsole.
\item
  \textbf{CSV lässt sich nicht laden} → Datei lokal in \texttt{data/}
  speichern und von dort lesen (CORS-Beschränkungen). Die
  OWID-Grapher-URLs ändern sich gelegentlich; im Zweifel den
  CSV-Download direkt von der Grapher-Seite holen.
\item
  \textbf{Diagramm zu klein} → \texttt{fig-width:\ 8} und
  \texttt{fig-height:\ 5} im Chunk-Header setzen.
\item
  \textbf{Spaltenname unbekannt} → \texttt{colnames(owid)} in der
  R-Konsole laufen lassen, um die exakte Schreibweise zu sehen (manche
  Spalten haben Sonderzeichen).
\end{itemize}

\end{tcolorbox}

\section{Tutor}\label{tutor-4}

\begin{tcolorbox}[enhanced jigsaw, opacitybacktitle=0.6, leftrule=.75mm, title=\textcolor{quarto-callout-tip-color}{\faLightbulb}\hspace{0.5em}{Tutor öffnen}, colframe=quarto-callout-tip-color-frame, arc=.35mm, toprule=.15mm, colbacktitle=quarto-callout-tip-color!10!white, bottomrule=.15mm, colback=white, rightrule=.15mm, breakable, toptitle=1mm, bottomtitle=1mm, opacityback=0, titlerule=0mm, coltitle=black, left=2mm]

\href{https://chatgpt.com/g/g-69fec93348108191af1ffb8ce0cf6712-tutor-interaktive-skripte}{Tutor
„Interaktive Skripte'' →}

Vorschlag-Prompt für dieses Kapitel:

\begin{quote}
Ich will ein interaktives Diagramm in mein Quarto-Skript einbauen.
Daten: {[}URL oder lokale Datei{]}. Frage: {[}welche Variablen, welcher
Vergleich{]}. Ich bin R-Anfänger. Schreibe mir einen kommentierten
Code-Chunk, in dem jede Zeile mit einer Lay-Erklärung versehen ist.
Falls die Daten nicht ladbar sind, schlag einen Workaround vor.
\end{quote}

\end{tcolorbox}

\chapter*{Literatur}\label{literatur-4}
\addcontentsline{toc}{chapter}{Literatur}

\markboth{Literatur}{Literatur}

\phantomsection\label{refs}
\begin{CSLReferences}{1}{0}
\bibitem[\citeproctext]{ref-allaire2024}
Allaire, J. J., Teague, C., Scheidegger, C., Xie, Y., \& Dervieux, C.
(2024). \emph{Quarto: An Open-Source Scientific and Technical Publishing
System}. \url{https://quarto.org/}

\bibitem[\citeproctext]{ref-bergstrom2024}
Bergstrom, C. T., \& West, J. D. (2024). \emph{Calling Bullshit:
Wissenschaftliches Arbeiten und kritisches Denken}.
\url{https://publish.obsidian.md/wagp-calling-bullshit/}

\bibitem[\citeproctext]{ref-bryan2018}
Bryan, J. (2018). Excuse Me, Do You Have a Moment to Talk About Version
Control? \emph{The American Statistician}, \emph{72}(1), 20--27.
\url{https://doi.org/10.1080/00031305.2017.1399928}

\bibitem[\citeproctext]{ref-fiorella2016}
Fiorella, L., \& Mayer, R. E. (2016). Eight Ways to Promote Generative
Learning. \emph{Educational Psychology Review}, \emph{28}(4), 717--741.
\url{https://doi.org/10.1007/s10648-015-9348-9}

\bibitem[\citeproctext]{ref-healy2018}
Healy, K. (2018). \emph{Data Visualization: A Practical Introduction}.
Princeton University Press.

\bibitem[\citeproctext]{ref-karpicke2011}
Karpicke, J. D., \& Blunt, J. R. (2011). Retrieval Practice Produces
More Learning than Elaborative Studying with Concept Mapping.
\emph{Science}, \emph{331}(6018), 772--775.
\url{https://doi.org/10.1126/science.1199327}

\bibitem[\citeproctext]{ref-kasneci2023}
Kasneci, E., Sessler, K., Küchemann, S., Bannert, M., Dementieva, D.,
Fischer, F., Gasser, U., Groh, G., Günnemann, S., Hüllermeier, E.,
Krusche, S., Kutyniok, G., Michaeli, T., Nerdel, C., Pfeffer, J.,
Poquet, O., Sailer, M., Schmidt, A., Seidel, T., \ldots{} Kasneci, G.
(2023). {ChatGPT} for Good? On Opportunities and Challenges of Large
Language Models for Education. \emph{Learning and Individual
Differences}, \emph{103}, 102274.
\url{https://doi.org/10.1016/j.lindif.2023.102274}

\bibitem[\citeproctext]{ref-lakens2022}
Lakens, D. (2022). \emph{Improving Your Statistical Inferences}.
\url{https://doi.org/10.5281/zenodo.6409077}

\bibitem[\citeproctext]{ref-mayer2021}
Mayer, R. E. (2021). \emph{Multimedia Learning} (3. Aufl.). Cambridge
University Press. \url{https://doi.org/10.1017/9781316941355}

\bibitem[\citeproctext]{ref-mollick2023}
Mollick, E. R., \& Mollick, L. (2023). Assigning {AI}: Seven Approaches
for Students, with Prompts. \emph{SSRN Working Paper}.
\url{https://doi.org/10.2139/ssrn.4475995}

\bibitem[\citeproctext]{ref-roediger2006}
Roediger, H. L., \& Karpicke, J. D. (2006). Test-Enhanced Learning:
Taking Memory Tests Improves Long-Term Retention. \emph{Psychological
Science}, \emph{17}(3), 249--255.
\url{https://doi.org/10.1111/j.1467-9280.2006.01693.x}

\bibitem[\citeproctext]{ref-owid2024}
Roser, M., Ritchie, H., Mathieu, E., \& Ortiz-Ospina, E. (2024).
\emph{Our World in Data}. \url{https://ourworldindata.org/}

\bibitem[\citeproctext]{ref-tufte2001}
Tufte, E. R. (2001). \emph{The Visual Display of Quantitative
Information} (2. Aufl.). Graphics Press.

\bibitem[\citeproctext]{ref-wickham2019}
Wickham, H., Averick, M., Bryan, J., Chang, W., McGowan, L. D.,
François, R., Grolemund, G., Hayes, A., Henry, L., Hester, J., Kuhn, M.,
Pedersen, T. L., Miller, E., Bache, S. M., Müller, K., Ooms, J.,
Robinson, D., Seidel, D. P., Spinu, V., \ldots{} Yutani, H. (2019).
Welcome to the Tidyverse. \emph{Journal of Open Source Software},
\emph{4}(43), 1686. \url{https://doi.org/10.21105/joss.01686}

\end{CSLReferences}

\chapter{5 --- iFrames, Spiele und
KI-Tutor}\label{iframes-spiele-und-ki-tutor}

\subsubsection{Was du nach diesem Kapitel
kannst}\label{was-du-nach-diesem-kapitel-kannst-5}

\begin{itemize}
\tightlist
\item
  ein \textbf{interaktives Quiz oder Spiel} als iFrame in dein Skript
  einbinden,
\item
  den Unterschied zwischen statischem HTML, \textbf{H5P} und externer
  Einbettung benennen,
\item
  einen \textbf{KI-Tutor} als Link einbauen, der die Lernenden begleitet
  --- didaktisch begründet.
\end{itemize}

\section{Worum es im Kern geht}\label{worum-es-im-kern-geht-4}

Aktivität schlägt Konsum. Lernende, die zwischendurch etwas tun ---
antworten, ziehen, ausprobieren --- behalten mehr und verstehen besser
(\citeproc{ref-fiorella2016}{Fiorella \& Mayer, 2016};
\citeproc{ref-karpicke2011}{Karpicke \& Blunt, 2011}). Interaktive
Elemente sind das Werkzeug, das diese Aktivität in einem statischen
Skript ermöglicht. Bei generativer KI als Tutor kommt eine zweite Chance
hinzu: ein adaptiver Gesprächspartner, der auf die jeweilige
Verständnisstufe reagiert (\citeproc{ref-kasneci2023}{Kasneci et al.,
2023}; \citeproc{ref-mollick2023}{Mollick \& Mollick, 2023}).

\section{iFrame in einer Zeile}\label{iframe-in-einer-zeile}

Ein iFrame bettet eine fremde Seite in deine ein:

\begin{Shaded}
\begin{Highlighting}[]
\DataTypeTok{\textless{}}\KeywordTok{iframe}\OtherTok{ src}\OperatorTok{=}\StringTok{"https://h5p.org/h5p/embed/617"}\OtherTok{ width}\OperatorTok{=}\StringTok{"100\%"}\OtherTok{ height}\OperatorTok{=}\StringTok{"500"}\OtherTok{ frameborder}\OperatorTok{=}\StringTok{"0"}\DataTypeTok{\textgreater{}\textless{}/}\KeywordTok{iframe}\DataTypeTok{\textgreater{}}
\end{Highlighting}
\end{Shaded}

Höhe nach Inhalt anpassen, Breite immer 100 \%, \texttt{frameborder="0"}
für ein sauberes Aussehen.

\section{Drei Wege, ein Spiel
einzubauen}\label{drei-wege-ein-spiel-einzubauen}

\begin{longtable}[]{@{}
  >{\raggedright\arraybackslash}p{(\linewidth - 6\tabcolsep) * \real{0.2500}}
  >{\raggedright\arraybackslash}p{(\linewidth - 6\tabcolsep) * \real{0.2500}}
  >{\raggedright\arraybackslash}p{(\linewidth - 6\tabcolsep) * \real{0.2500}}
  >{\raggedright\arraybackslash}p{(\linewidth - 6\tabcolsep) * \real{0.2500}}@{}}
\toprule\noalign{}
\begin{minipage}[b]{\linewidth}\raggedright
Variante
\end{minipage} & \begin{minipage}[b]{\linewidth}\raggedright
Aufwand
\end{minipage} & \begin{minipage}[b]{\linewidth}\raggedright
Vorteil
\end{minipage} & \begin{minipage}[b]{\linewidth}\raggedright
Nachteil
\end{minipage} \\
\midrule\noalign{}
\endhead
\bottomrule\noalign{}
\endlastfoot
\textbf{A --- Statisches HTML} (eigenes Widget) & hoch zur Erstellung,
niedrig danach & volle Kontrolle, kein Konto, läuft offline & erfordert
HTML/JS-Kenntnisse oder KI-Hilfe \\
\textbf{B --- H5P} (h5p.org / Moodle / WordPress) & mittel & viele
fertige Vorlagen (Quiz, Drag-Drop, Memory) & erfordert Konto, externe
Abhängigkeit \\
\textbf{C --- externer Editor} (Genially, Kahoot, Quizizz) & niedrig &
schnell, hübsch, kollaborativ & Werbung, externe Datenhaltung, oft Login
nötig \\
\end{longtable}

\begin{tcolorbox}[enhanced jigsaw, opacitybacktitle=0.6, leftrule=.75mm, title=\textcolor{quarto-callout-tip-color}{\faLightbulb}\hspace{0.5em}{Mach mit (Variante A --- statisches HTML)}, colframe=quarto-callout-tip-color-frame, arc=.35mm, toprule=.15mm, colbacktitle=quarto-callout-tip-color!10!white, bottomrule=.15mm, colback=white, rightrule=.15mm, breakable, toptitle=1mm, bottomtitle=1mm, opacityback=0, titlerule=0mm, coltitle=black, left=2mm]

Schau dir das \textbf{Lernspiel aus Kapitel 1}
(\href{../interactions/lernspiel-kapitel1.html}{interactions/lernspiel-kapitel1.html})
im Browser an. Öffne die Datei in RStudio (im \textbf{Files}-Tab
anklicken → Open in editor) --- sie ist eine einzige HTML-Datei, ohne
Backend. Diesen Bauplan kannst du selbst (oder dein KI-Tutor) für ein
eigenes Spiel im gewählten Strang nutzen.

\textbf{Strang-Beispiele:}

\begin{itemize}
\tightlist
\item
  Strang A --- Drag-and-Drop: „Welche Aussage ist eine Hypothese, welche
  eine Behauptung?''
\item
  Strang B --- Multiple Choice: „Welche der drei Quellen ist
  primärwissenschaftlich?''
\item
  Strang C --- Memory: „Methodenbegriff zu transparenter Beschreibung''
\end{itemize}

\end{tcolorbox}

\begin{tcolorbox}[enhanced jigsaw, opacitybacktitle=0.6, leftrule=.75mm, title=\textcolor{quarto-callout-tip-color}{\faLightbulb}\hspace{0.5em}{Mach mit (Variante B --- H5P, optional)}, colframe=quarto-callout-tip-color-frame, arc=.35mm, toprule=.15mm, colbacktitle=quarto-callout-tip-color!10!white, bottomrule=.15mm, colback=white, rightrule=.15mm, breakable, toptitle=1mm, bottomtitle=1mm, opacityback=0, titlerule=0mm, coltitle=black, left=2mm]

Lege ein kostenloses Konto auf \url{https://h5p.org} an. Wähle
\textbf{Drag the Words}, baue ein Mini-Quiz zu drei Begriffen aus deinem
Strang. Klicke „Embed'', kopiere den iFrame-Code, füge ihn in dein
Kapitel ein. Render.

\end{tcolorbox}

\section{KI-Tutor als Link}\label{ki-tutor-als-link}

Der Tutor läuft \textbf{nicht} im iFrame, sondern in einem zweiten Tab.
Du legst einen GPT (oder Custom GPT) mit System-Prompt an, holst dir den
öffentlichen Link, baust ihn als Markdown-Link an passende Stellen im
Skript ein. Drei Hosting-Varianten zur Auswahl:

\subsection{Option 1 --- Custom GPT auf ChatGPT (heute
genutzt)}\label{option-1-custom-gpt-auf-chatgpt-heute-genutzt}

Einfachste Variante, weltweit verfügbar, kein Hochschul-Login nötig. Im
Skript einbauen:

\begin{Shaded}
\begin{Highlighting}[]
\CommentTok{[}\OtherTok{Frag den Skript{-}Tutor →}\CommentTok{](https://chatgpt.com/g/g{-}69fec93348108191af1ffb8ce0cf6712{-}tutor{-}interaktive{-}skripte)}\NormalTok{\{target="\_blank"\}}
\end{Highlighting}
\end{Shaded}

Der vorbereitete Tutor für diesen Workshop liegt unter dem oben
verlinkten URL --- der System-Prompt steckt dort schon drin. Die
Lernenden klicken, ChatGPT-Login öffnet sich (kostenlos), der Tutor
antwortet sofort im Workshop-Kontext.

\begin{tcolorbox}[enhanced jigsaw, opacitybacktitle=0.6, leftrule=.75mm, title=\textcolor{quarto-callout-note-color}{\faInfo}\hspace{0.5em}{Hinweis}, colframe=quarto-callout-note-color-frame, arc=.35mm, toprule=.15mm, colbacktitle=quarto-callout-note-color!10!white, bottomrule=.15mm, colback=white, rightrule=.15mm, breakable, toptitle=1mm, bottomtitle=1mm, opacityback=0, titlerule=0mm, coltitle=black, left=2mm]

\textbf{Datenschutz:} ChatGPT speichert Konversationen auf US-Servern.
Für offene Lehre ohne personenbezogene Daten unkritisch. Wenn du
europäische Datenhaltung brauchst, nimm Option 2.

\end{tcolorbox}

\subsection{Option 2 --- Custom GPT in der
AcademicCloud}\label{option-2-custom-gpt-in-der-academiccloud}

Variante mit europäischer Datenhaltung und Hochschul-Login. Du legst
einen GPT mit dem vorbereiteten System-Prompt an (siehe
\href{../appendix/tutor-prompt.qmd}{Anhang Tutor-Prompt}) und verlinkst
ihn:

\begin{Shaded}
\begin{Highlighting}[]
\CommentTok{[}\OtherTok{Frag den Skript{-}Tutor →}\CommentTok{](https://chat.academiccloud.de/chat/DEIN{-}LINK)}\NormalTok{\{target="\_blank"\}}
\end{Highlighting}
\end{Shaded}

Lernende loggen sich mit Hochschul-Credentials ein.

\subsection{Option 3 --- eigenes statisches Widget oder voll gehosteter
Tutor}\label{option-3-eigenes-statisches-widget-oder-voll-gehosteter-tutor}

Für Fortgeschrittene: ein HTML-Widget, das eine OpenAI-kompatible API
direkt anspricht (Vorbild: \texttt{\_relevance-widget.qmd} in
\url{https://github.com/TH-Koln-Bartnik/genai4teaching}), oder eine
eigene Web-App auf AWS Amplify mit RAG über Skriptinhalte (Vorbild:
\href{https://main.difctumi6okdj.amplifyapp.com/}{BWL1 Tutor-Chatbot}).
Mollick und Mollick (\citeproc{ref-mollick2023}{2023}) geben einen
Überblick über die didaktischen Designentscheidungen.

\begin{tcolorbox}[enhanced jigsaw, opacitybacktitle=0.6, leftrule=.75mm, title=\textcolor{quarto-callout-tip-color}{\faLightbulb}\hspace{0.5em}{Mach mit}, colframe=quarto-callout-tip-color-frame, arc=.35mm, toprule=.15mm, colbacktitle=quarto-callout-tip-color!10!white, bottomrule=.15mm, colback=white, rightrule=.15mm, breakable, toptitle=1mm, bottomtitle=1mm, opacityback=0, titlerule=0mm, coltitle=black, left=2mm]

Wir nutzen heute den \textbf{vorbereiteten Workshop-Tutor} auf ChatGPT
--- du musst keinen eigenen GPT bauen, sondern nur den Link einbauen:

\begin{Shaded}
\begin{Highlighting}[]
\CommentTok{[}\OtherTok{Frag den Skript{-}Tutor →}\CommentTok{](https://chatgpt.com/g/g{-}69fec93348108191af1ffb8ce0cf6712{-}tutor{-}interaktive{-}skripte)}\NormalTok{\{target="\_blank"\}}
\end{Highlighting}
\end{Shaded}

Setze diesen Link an passende Stelle in dein Skript (z. B. in einen
Callout-Tip am Kapitelende). Render und prüf den Klick.

Wer einen \textbf{eigenen} Tutor mit eigenem System-Prompt bauen will,
springt zum \href{../appendix/tutor-prompt.qmd}{Tutor-Prompt-Anhang} und
legt sich Option 1 oder Option 2 selbst an.

\end{tcolorbox}

\section{\texorpdfstring{Wann \textbf{kein}
KI-Tutor?}{Wann kein KI-Tutor?}}\label{wann-kein-ki-tutor}

Drei Situationen, in denen du es lassen solltest:

\begin{itemize}
\tightlist
\item
  \textbf{Klar definierte Faktenfragen} --- wenn die Antwort eindeutig
  ist und in einem Lehrbuch steht, ist ein Quiz besser.
\item
  \textbf{Sensible persönliche Themen} --- der Tutor soll Lernen
  begleiten, keine Beratung ersetzen.
\item
  \textbf{Zeitkritisches Pflichtwissen} --- wenn die Lernenden in zehn
  Minuten etwas können müssen, ist Frontalanleitung effizienter als ein
  Gespräch (\citeproc{ref-kasneci2023}{Kasneci et al., 2023}).
\end{itemize}

\begin{tcolorbox}[enhanced jigsaw, opacitybacktitle=0.6, leftrule=.75mm, title=\textcolor{quarto-callout-warning-color}{\faExclamationTriangle}\hspace{0.5em}{Was, wenn \ldots{}}, colframe=quarto-callout-warning-color-frame, arc=.35mm, toprule=.15mm, colbacktitle=quarto-callout-warning-color!10!white, bottomrule=.15mm, colback=white, rightrule=.15mm, breakable, toptitle=1mm, bottomtitle=1mm, opacityback=0, titlerule=0mm, coltitle=black, left=2mm]

\begin{itemize}
\tightlist
\item
  \textbf{iFrame zeigt nichts an?} → URL prüfen, oft fehlt
  \texttt{https://}. Manche Seiten erlauben kein Embedding (HTTP-Header
  \texttt{X-Frame-Options}) --- dann nur als Link verwenden.
\item
  \textbf{iFrame zu kurz, Inhalt abgeschnitten?} → \texttt{height} höher
  setzen, oder \texttt{style="height:\ 90vh"} für 90 \% der
  Bildschirmhöhe.
\item
  \textbf{AcademicCloud-Link funktioniert für andere nicht?} → Sharing
  aktivieren in den GPT-Einstellungen, „öffentlich'' oder „mit Link''.
\item
  \textbf{Tutor antwortet auf alles, auch außerhalb des Themas?} →
  System-Prompt nachschärfen, explizit Grenzen setzen („Antworte nur zu
  Quarto-Workflow-Fragen'').
\end{itemize}

\end{tcolorbox}

\section{Tutor}\label{tutor-5}

\begin{tcolorbox}[enhanced jigsaw, opacitybacktitle=0.6, leftrule=.75mm, title=\textcolor{quarto-callout-tip-color}{\faLightbulb}\hspace{0.5em}{Tutor öffnen}, colframe=quarto-callout-tip-color-frame, arc=.35mm, toprule=.15mm, colbacktitle=quarto-callout-tip-color!10!white, bottomrule=.15mm, colback=white, rightrule=.15mm, breakable, toptitle=1mm, bottomtitle=1mm, opacityback=0, titlerule=0mm, coltitle=black, left=2mm]

\href{https://chatgpt.com/g/g-69fec93348108191af1ffb8ce0cf6712-tutor-interaktive-skripte}{Tutor
„Interaktive Skripte'' →}

Vorschlag-Prompt für dieses Kapitel:

\begin{quote}
Ich will {[}ein iFrame-Quiz / einen KI-Tutor-Link / ein eigenes
HTML-Widget{]} in mein Quarto-Skript einbauen. Mein Stand: {[}\ldots{]}.
Schlag mir den einfachsten ersten Schritt vor und schreib den
Code-Schnipsel dazu. Erkläre vorher in zwei Sätzen, was ich da
eigentlich tue.
\end{quote}

\end{tcolorbox}

\chapter*{Literatur}\label{literatur-5}
\addcontentsline{toc}{chapter}{Literatur}

\markboth{Literatur}{Literatur}

\phantomsection\label{refs}
\begin{CSLReferences}{1}{0}
\bibitem[\citeproctext]{ref-allaire2024}
Allaire, J. J., Teague, C., Scheidegger, C., Xie, Y., \& Dervieux, C.
(2024). \emph{Quarto: An Open-Source Scientific and Technical Publishing
System}. \url{https://quarto.org/}

\bibitem[\citeproctext]{ref-bergstrom2024}
Bergstrom, C. T., \& West, J. D. (2024). \emph{Calling Bullshit:
Wissenschaftliches Arbeiten und kritisches Denken}.
\url{https://publish.obsidian.md/wagp-calling-bullshit/}

\bibitem[\citeproctext]{ref-bryan2018}
Bryan, J. (2018). Excuse Me, Do You Have a Moment to Talk About Version
Control? \emph{The American Statistician}, \emph{72}(1), 20--27.
\url{https://doi.org/10.1080/00031305.2017.1399928}

\bibitem[\citeproctext]{ref-fiorella2016}
Fiorella, L., \& Mayer, R. E. (2016). Eight Ways to Promote Generative
Learning. \emph{Educational Psychology Review}, \emph{28}(4), 717--741.
\url{https://doi.org/10.1007/s10648-015-9348-9}

\bibitem[\citeproctext]{ref-healy2018}
Healy, K. (2018). \emph{Data Visualization: A Practical Introduction}.
Princeton University Press.

\bibitem[\citeproctext]{ref-karpicke2011}
Karpicke, J. D., \& Blunt, J. R. (2011). Retrieval Practice Produces
More Learning than Elaborative Studying with Concept Mapping.
\emph{Science}, \emph{331}(6018), 772--775.
\url{https://doi.org/10.1126/science.1199327}

\bibitem[\citeproctext]{ref-kasneci2023}
Kasneci, E., Sessler, K., Küchemann, S., Bannert, M., Dementieva, D.,
Fischer, F., Gasser, U., Groh, G., Günnemann, S., Hüllermeier, E.,
Krusche, S., Kutyniok, G., Michaeli, T., Nerdel, C., Pfeffer, J.,
Poquet, O., Sailer, M., Schmidt, A., Seidel, T., \ldots{} Kasneci, G.
(2023). {ChatGPT} for Good? On Opportunities and Challenges of Large
Language Models for Education. \emph{Learning and Individual
Differences}, \emph{103}, 102274.
\url{https://doi.org/10.1016/j.lindif.2023.102274}

\bibitem[\citeproctext]{ref-lakens2022}
Lakens, D. (2022). \emph{Improving Your Statistical Inferences}.
\url{https://doi.org/10.5281/zenodo.6409077}

\bibitem[\citeproctext]{ref-mayer2021}
Mayer, R. E. (2021). \emph{Multimedia Learning} (3. Aufl.). Cambridge
University Press. \url{https://doi.org/10.1017/9781316941355}

\bibitem[\citeproctext]{ref-mollick2023}
Mollick, E. R., \& Mollick, L. (2023). Assigning {AI}: Seven Approaches
for Students, with Prompts. \emph{SSRN Working Paper}.
\url{https://doi.org/10.2139/ssrn.4475995}

\bibitem[\citeproctext]{ref-roediger2006}
Roediger, H. L., \& Karpicke, J. D. (2006). Test-Enhanced Learning:
Taking Memory Tests Improves Long-Term Retention. \emph{Psychological
Science}, \emph{17}(3), 249--255.
\url{https://doi.org/10.1111/j.1467-9280.2006.01693.x}

\bibitem[\citeproctext]{ref-owid2024}
Roser, M., Ritchie, H., Mathieu, E., \& Ortiz-Ospina, E. (2024).
\emph{Our World in Data}. \url{https://ourworldindata.org/}

\bibitem[\citeproctext]{ref-tufte2001}
Tufte, E. R. (2001). \emph{The Visual Display of Quantitative
Information} (2. Aufl.). Graphics Press.

\bibitem[\citeproctext]{ref-wickham2019}
Wickham, H., Averick, M., Bryan, J., Chang, W., McGowan, L. D.,
François, R., Grolemund, G., Hayes, A., Henry, L., Hester, J., Kuhn, M.,
Pedersen, T. L., Miller, E., Bache, S. M., Müller, K., Ooms, J.,
Robinson, D., Seidel, D. P., Spinu, V., \ldots{} Yutani, H. (2019).
Welcome to the Tidyverse. \emph{Journal of Open Source Software},
\emph{4}(43), 1686. \url{https://doi.org/10.21105/joss.01686}

\end{CSLReferences}

\chapter{6 --- Slides mit Quarto
(optional)}\label{slides-mit-quarto-optional}

\begin{tcolorbox}[enhanced jigsaw, opacitybacktitle=0.6, leftrule=.75mm, title=\textcolor{quarto-callout-note-color}{\faInfo}\hspace{0.5em}{Optional, für Fortgeschrittene}, colframe=quarto-callout-note-color-frame, arc=.35mm, toprule=.15mm, colbacktitle=quarto-callout-note-color!10!white, bottomrule=.15mm, colback=white, rightrule=.15mm, breakable, toptitle=1mm, bottomtitle=1mm, opacityback=0, titlerule=0mm, coltitle=black, left=2mm]

Dieses Kapitel ist nicht Teil der gemeinsamen Übung. Wer mit Kapitel
1--5 früher fertig ist, kann hier weiterbauen. Vorlage und Inspiration:
Békés (2025),
\href{https://gabors-data-analysis.com/courses/da-w-ai-2025/da-w-ai-01-llm-course.html}{Course
slides for ``Data Analysis with AI''}.

\end{tcolorbox}

\subsubsection{Was du nach diesem Kapitel
kannst}\label{was-du-nach-diesem-kapitel-kannst-6}

\begin{itemize}
\tightlist
\item
  aus einem \texttt{.qmd}-Kapitel einen \textbf{Foliensatz} mit Quarto
  Reveal.js erzeugen,
\item
  die \textbf{TH-Köln-Farben} auf Slides übertragen,
\item
  den Foliensatz neben das Skript auf GitHub Pages stellen.
\end{itemize}

\section{Schnellbau in vier
Schritten}\label{schnellbau-in-vier-schritten}

\begin{enumerate}
\def\labelenumi{\arabic{enumi}.}
\tightlist
\item
  Neue Datei \texttt{slides/kapitel-1-slides.qmd} anlegen.
\item
  YAML-Header:
\end{enumerate}

\begin{Shaded}
\begin{Highlighting}[]
\PreprocessorTok{{-}{-}{-}}
\FunctionTok{title}\KeywordTok{:}\AttributeTok{ }\StringTok{"Kapitel 1 — Didaktik"}
\FunctionTok{author}\KeywordTok{:}\AttributeTok{ }\StringTok{"Dein Name"}
\FunctionTok{format}\KeywordTok{:}
\AttributeTok{  }\FunctionTok{revealjs}\KeywordTok{:}
\AttributeTok{    }\FunctionTok{theme}\KeywordTok{:}\AttributeTok{ }\KeywordTok{[}\AttributeTok{default}\KeywordTok{,}\AttributeTok{ ../styles.scss}\KeywordTok{]}
\AttributeTok{    }\FunctionTok{slide{-}number}\KeywordTok{:}\AttributeTok{ }\CharTok{true}
\AttributeTok{    }\FunctionTok{chalkboard}\KeywordTok{:}\AttributeTok{ }\CharTok{true}
\AttributeTok{    }\FunctionTok{incremental}\KeywordTok{:}\AttributeTok{ }\CharTok{true}
\AttributeTok{    }\FunctionTok{logo}\KeywordTok{:}\AttributeTok{ ../images/th{-}koeln{-}logo.png}
\PreprocessorTok{{-}{-}{-}}
\end{Highlighting}
\end{Shaded}

\begin{enumerate}
\def\labelenumi{\arabic{enumi}.}
\setcounter{enumi}{2}
\tightlist
\item
  Inhalt: \textbf{eine H2-Überschrift = eine Folie}, dazwischen
  Stichpunkte und Bilder. Speaker Notes mit \texttt{:::\ notes}.
\item
  \texttt{quarto\ render\ slides/kapitel-1-slides.qmd} → Slide-Deck als
  HTML.
\end{enumerate}

\section{TH-Köln-Brand auf Slides}\label{th-kuxf6ln-brand-auf-slides}

Die \texttt{\_brand.yml} aus deinem Skript wird auch von Reveal.js
gelesen. Zusätzlich einbinden:

\begin{Shaded}
\begin{Highlighting}[]
\FunctionTok{format}\KeywordTok{:}
\AttributeTok{  }\FunctionTok{revealjs}\KeywordTok{:}
\AttributeTok{    }\FunctionTok{brand}\KeywordTok{:}\AttributeTok{ ../\_brand.yml}
\AttributeTok{    }\FunctionTok{theme}\KeywordTok{:}\AttributeTok{ }\KeywordTok{[}\AttributeTok{default}\KeywordTok{,}\AttributeTok{ ../styles.scss}\KeywordTok{]}
\end{Highlighting}
\end{Shaded}

\section{Tutor}\label{tutor-6}

\begin{tcolorbox}[enhanced jigsaw, opacitybacktitle=0.6, leftrule=.75mm, title=\textcolor{quarto-callout-tip-color}{\faLightbulb}\hspace{0.5em}{Tutor öffnen}, colframe=quarto-callout-tip-color-frame, arc=.35mm, toprule=.15mm, colbacktitle=quarto-callout-tip-color!10!white, bottomrule=.15mm, colback=white, rightrule=.15mm, breakable, toptitle=1mm, bottomtitle=1mm, opacityback=0, titlerule=0mm, coltitle=black, left=2mm]

\href{https://chatgpt.com/g/g-69fec93348108191af1ffb8ce0cf6712-tutor-interaktive-skripte}{Tutor
„Interaktive Skripte'' →}

Vorschlag-Prompt für dieses Kapitel:

\begin{quote}
Ich will aus meinem bestehenden Quarto-Kapitel einen Foliensatz mit
Reveal.js machen. TH-Köln-Brand soll erhalten bleiben. Schlag mir die
einfachste Aufteilung vor: Welche Abschnitte werden eine Folie, welche
mehrere? Gib mir den YAML-Header.
\end{quote}

\end{tcolorbox}

\part{Teil 2 --- Quick-Starter}

\chapter{Quick-Starter Guide}\label{quick-starter-guide}

Von Null auf erstem Skript live im Netz --- alles aus RStudio

\hfill\break

\subsubsection{Was dieser Anhang
leistet}\label{was-dieser-anhang-leistet}

Eine \textbf{Schritt-für-Schritt-Anleitung für komplette Anfänger},
parallel zum Workshop nutzbar. Wer alles installiert hat und einmal
pushen will: 30 Minuten reine Klickarbeit. \textbf{RStudio ist dein
einziges Editor-Programm} --- Terminal, Files-Pane, Build-Pane und
Console liegen alle dort.

\section{Checkliste vor dem Workshop}\label{checkliste-vor-dem-workshop}

\begin{longtable}[]{@{}
  >{\raggedright\arraybackslash}p{(\linewidth - 6\tabcolsep) * \real{0.2500}}
  >{\raggedright\arraybackslash}p{(\linewidth - 6\tabcolsep) * \real{0.2500}}
  >{\raggedright\arraybackslash}p{(\linewidth - 6\tabcolsep) * \real{0.2500}}
  >{\raggedright\arraybackslash}p{(\linewidth - 6\tabcolsep) * \real{0.2500}}@{}}
\toprule\noalign{}
\begin{minipage}[b]{\linewidth}\raggedright
\end{minipage} & \begin{minipage}[b]{\linewidth}\raggedright
Schritt
\end{minipage} & \begin{minipage}[b]{\linewidth}\raggedright
Wo?
\end{minipage} & \begin{minipage}[b]{\linewidth}\raggedright
Zeit
\end{minipage} \\
\midrule\noalign{}
\endhead
\bottomrule\noalign{}
\endlastfoot
☐ & Quarto installieren & \url{https://quarto.org/docs/get-started/} & 5
min \\
☐ & R installieren & \url{https://cran.r-project.org/} & 5 min \\
☐ & RStudio Desktop (kostenlos) &
\url{https://posit.co/download/rstudio-desktop/} & 5 min \\
☐ & GitHub Desktop & \url{https://desktop.github.com/} & 5 min \\
☐ & GitHub-Konto erstellen & \url{https://github.com/signup} & 5 min \\
☐ & Zotero + Better BibTeX & \url{https://www.zotero.org/} & 10 min \\
☐ & ChatGPT-Konto und Tutor-Link testen &
\url{https://chatgpt.com/g/g-69fec93348108191af1ffb8ce0cf6712-tutor-interaktive-skripte}
& 2 min \\
\end{longtable}

Insgesamt rund 40 Minuten. Wer schon hat, was er braucht, springt
weiter.

\section{RStudio-Cockpit auf einen
Blick}\label{rstudio-cockpit-auf-einen-blick}

\begin{longtable}[]{@{}
  >{\raggedright\arraybackslash}p{(\linewidth - 4\tabcolsep) * \real{0.3333}}
  >{\raggedright\arraybackslash}p{(\linewidth - 4\tabcolsep) * \real{0.3333}}
  >{\raggedright\arraybackslash}p{(\linewidth - 4\tabcolsep) * \real{0.3333}}@{}}
\toprule\noalign{}
\begin{minipage}[b]{\linewidth}\raggedright
Bereich
\end{minipage} & \begin{minipage}[b]{\linewidth}\raggedright
Position
\end{minipage} & \begin{minipage}[b]{\linewidth}\raggedright
Wofür im Workshop?
\end{minipage} \\
\midrule\noalign{}
\endhead
\bottomrule\noalign{}
\endlastfoot
\textbf{Editor} & oben links & \texttt{.qmd}-Dateien schreiben \\
\textbf{Console} & unten links, Tab 1 & R-Pakete installieren
(\texttt{install.packages(...)}) \\
\textbf{Terminal} & unten links, Tab 2 & \texttt{quarto\ render},
\texttt{quarto\ preview} \\
\textbf{Files} & unten rechts, Tab 1 & Ordner und Dateien anlegen,
Dateien öffnen \\
\textbf{Build} & oben rechts, Tab 5 & Knopf „Render Book'' als
Klick-Alternative \\
\end{longtable}

Drei Shortcuts, die du dir merken solltest: \textbf{Strg+Shift+K}
rendert die aktuelle Datei, \textbf{Strg+Shift+B} rendert das ganze
Buch, \textbf{Alt+Shift+R} öffnet einen neuen Terminal-Tab.

\section{Schritt-für-Schritt: dein erstes
Skript}\label{schritt-fuxfcr-schritt-dein-erstes-skript}

\subsection{1. Neues Quarto-Buch
anlegen}\label{neues-quarto-buch-anlegen}

In RStudio: \textbf{File → New Project → New Directory → Quarto Book →
Create Project}.

\subsection{2. Erste Vorschau}\label{erste-vorschau}

Im \textbf{Terminal-Tab} (unten links, neben Console):

\begin{Shaded}
\begin{Highlighting}[]
\ExtensionTok{quarto}\NormalTok{ preview}
\end{Highlighting}
\end{Shaded}

Browser öffnet sich. Du siehst die Beispielinhalte. \textbf{Strg+C} im
Terminal beendet den Preview-Server.

\subsection{\texorpdfstring{3. \texttt{\_quarto.yml} minimal
anpassen}{3. \_quarto.yml minimal anpassen}}\label{quarto.yml-minimal-anpassen}

Im \textbf{Editor} öffnest du \texttt{\_quarto.yml} und ersetzt den
Inhalt durch:

\begin{Shaded}
\begin{Highlighting}[]
\FunctionTok{project}\KeywordTok{:}
\AttributeTok{  }\FunctionTok{type}\KeywordTok{:}\AttributeTok{ book}
\AttributeTok{  }\FunctionTok{output{-}dir}\KeywordTok{:}\AttributeTok{ docs}

\FunctionTok{book}\KeywordTok{:}
\AttributeTok{  }\FunctionTok{title}\KeywordTok{:}\AttributeTok{ }\StringTok{"Mein erstes Skript"}
\AttributeTok{  }\FunctionTok{chapters}\KeywordTok{:}
\AttributeTok{    }\KeywordTok{{-}}\AttributeTok{ index.qmd}
\AttributeTok{    }\KeywordTok{{-}}\AttributeTok{ intro.qmd}

\FunctionTok{format}\KeywordTok{:}
\AttributeTok{  }\FunctionTok{html}\KeywordTok{:}
\AttributeTok{    }\FunctionTok{theme}\KeywordTok{:}\AttributeTok{ cosmo}
\AttributeTok{    }\FunctionTok{toc}\KeywordTok{:}\AttributeTok{ }\CharTok{true}
\end{Highlighting}
\end{Shaded}

\texttt{output-dir:\ docs} ist wichtig --- GitHub Pages erwartet das
später.

\subsection{4. Ordner für Bilder, Daten und Widgets
anlegen}\label{ordner-fuxfcr-bilder-daten-und-widgets-anlegen}

Im \textbf{Files-Pane} (unten rechts) klickst du auf \textbf{New Folder}
und legst nacheinander an: \texttt{images/}, \texttt{interactions/},
\texttt{include/}, \texttt{data/}. So weißt du immer, wo etwas
hingehört.

\subsection{5. Erstes Kapitel schreiben}\label{erstes-kapitel-schreiben}

Öffne \texttt{index.qmd} im Editor, ersetze den Beispielinhalt durch:

\begin{Shaded}
\begin{Highlighting}[]
\FunctionTok{\# Willkommen}

\NormalTok{Mein erstes Skript zum Thema }\CommentTok{[}\OtherTok{Strang A/B/C}\CommentTok{]}\NormalTok{.}

\FunctionTok{\#\# Was ich hier behandele}

\SpecialStringTok{{-} }\NormalTok{Punkt 1}
\SpecialStringTok{{-} }\NormalTok{Punkt 2}
\SpecialStringTok{{-} }\NormalTok{Punkt 3}
\end{Highlighting}
\end{Shaded}

\subsection{6. Render}\label{render}

Im Terminal:

\begin{Shaded}
\begin{Highlighting}[]
\ExtensionTok{quarto}\NormalTok{ render}
\end{Highlighting}
\end{Shaded}

Oder Knopf \textbf{Build → Render Book} klicken. Im
\texttt{docs/}-Ordner liegt jetzt die fertige Webseite.

\subsection{7. GitHub Desktop: Repo
anlegen}\label{github-desktop-repo-anlegen}

GitHub Desktop → \textbf{File → Add Local Repository} → Projektordner
wählen → \textbf{Create a repository} zustimmen → Name:
\texttt{mein-erstes-skript} → \textbf{Publish repository} (Public).

\subsection{8. GitHub Pages aktivieren}\label{github-pages-aktivieren}

Auf \url{https://github.com} im Repo: \textbf{Settings → Pages → Source:
Deploy from a branch → Branch: main, Folder: /docs → Save}.

Nach 1--2 Minuten ist dein Skript unter
\texttt{https://DEIN-USERNAME.github.io/mein-erstes-skript/} live.

\subsection{9. Erste Änderung
publizieren}\label{erste-uxe4nderung-publizieren}

Bearbeite \texttt{index.qmd}, render lokal (Terminal:
\texttt{quarto\ render} oder Knopf), in GitHub Desktop \textbf{Commit →
Push origin}. Webseite aktualisiert sich automatisch in 30--60 Sekunden.

\section{Was als Nächstes?}\label{was-als-nuxe4chstes}

Springe zurück in \hyperref[a-inhaltsverzeichnis-und-literatur]{Kapitel
3a --- Übersicht} und baue Inhaltsverzeichnis und Literaturverwaltung
ein. Dann zu \hyperref[b-th-kuxf6ln-stil-einbauen]{Kapitel 3b --- Stil}.

\chapter{Trouble-Shooting}\label{trouble-shooting}

Die zehn häufigsten Probleme und wie du sie löst

\hfill\break

Jeder Punkt: \textbf{Symptom → Ursache → Lösung}. Wenn nichts hilft:
Tutor fragen mit dem Prompt am Ende.

\section{1. YAML-Fehler beim Rendern}\label{yaml-fehler-beim-rendern}

\textbf{Symptom:} \texttt{Error:\ Failed\ to\ load\ YAML\ metadata}.
\textbf{Ursache:} Einrückung mit Tabs statt Leerzeichen, oder fehlender
Doppelpunkt. \textbf{Lösung:} YAML-Block prüfen, immer \textbf{zwei
Leerzeichen} einrücken, keine Tabs. Online-Validator:
\url{https://www.yamllint.com}.

\section{2. GitHub Pages zeigt 404}\label{github-pages-zeigt-404}

\textbf{Symptom:} Dein Skript ist online, aber die URL liefert „404 ---
File not found''. \textbf{Ursache:} \texttt{output-dir:\ docs} fehlt in
\texttt{\_quarto.yml}, oder der Pages-Branch ist falsch eingestellt.
\textbf{Lösung:} In \texttt{\_quarto.yml} \texttt{output-dir:\ docs}
setzen, render, push. In \textbf{Settings → Pages}: Branch
\texttt{main}, Folder \texttt{/docs}.

\section{\texorpdfstring{3. \texttt{docs/} wird nicht
committet}{3. docs/ wird nicht committet}}\label{docs-wird-nicht-committet}

\textbf{Symptom:} Nach Push fehlt der \texttt{docs/}-Ordner auf GitHub.
\textbf{Ursache:} \texttt{.gitignore} schließt \texttt{/docs/} aus.
\textbf{Lösung:} Zeile \texttt{/docs/} aus \texttt{.gitignore}
entfernen, dann commit und push.

\section{\texorpdfstring{4. Zotero-Zitat erscheint als
\texttt{{[}?{]}}}{4. Zotero-Zitat erscheint als {[}?{]}}}\label{zotero-zitat-erscheint-als}

\textbf{Symptom:} Zitate werden mit Fragezeichen statt Autor und Jahr
gerendert. \textbf{Ursache:} Citation Key in der \texttt{.bib}-Datei
stimmt nicht mit dem Key im Text überein. \textbf{Lösung:}
\texttt{references.bib} öffnen, Schlüssel nachsehen, im Text exakt so
verwenden. Better BibTeX in Zotero exportiert stabile Keys (z. B.
\texttt{mayer2021}).

\section{5. iFrame ist leer}\label{iframe-ist-leer}

\textbf{Symptom:} Der iFrame-Bereich ist sichtbar, aber leer oder zeigt
eine Fehlermeldung. \textbf{Ursache:} Die externe Seite verbietet
Einbettung (HTTP-Header \texttt{X-Frame-Options}). \textbf{Lösung:}
Statt iFrame einen normalen Markdown-Link verwenden, oder eine andere
Quelle wählen.

\section{6. CSV lässt sich nicht
laden}\label{csv-luxe4sst-sich-nicht-laden}

\textbf{Symptom:} \texttt{read\_csv()} wirft \texttt{CORS\ error} oder
\texttt{Forbidden}. \textbf{Ursache:} Der Server der Datenquelle erlaubt
keinen Cross-Origin-Zugriff aus dem Browser. \textbf{Lösung:} Datei
lokal in \texttt{data/} speichern und von dort laden.

\section{7. R-Paket fehlt}\label{r-paket-fehlt}

\textbf{Symptom:}
\texttt{Error\ in\ library(plotly):\ there\ is\ no\ package\ called\ \textquotesingle{}plotly\textquotesingle{}}.
\textbf{Ursache:} Paket nicht installiert. \textbf{Lösung:} In der
R-Konsole \texttt{install.packages("plotly")}. Bei mehreren:
\texttt{install.packages(c("plotly",\ "dplyr",\ "readr"))}.

\section{8. Quarto-Version zu alt}\label{quarto-version-zu-alt}

\textbf{Symptom:} Neue Features (z. B. \texttt{\_brand.yml})
funktionieren nicht. \textbf{Ursache:} Lokale Quarto-Installation ist
älter als 1.5. \textbf{Lösung:} \texttt{quarto\ -\/-version} prüfen, von
\url{https://quarto.org/docs/get-started/} aktualisieren.

\section{9. GitHub Desktop fragt immer nach
Anmeldung}\label{github-desktop-fragt-immer-nach-anmeldung}

\textbf{Symptom:} Bei jedem Push wird das Passwort erneut verlangt.
\textbf{Ursache:} Personal Access Token läuft ab oder Kontostatus
inkonsistent. \textbf{Lösung:} In GitHub Desktop \textbf{File → Options
→ Accounts → Sign Out}, dann neu anmelden via Browser.

\section{10. Logo wird nicht angezeigt}\label{logo-wird-nicht-angezeigt}

\textbf{Symptom:} Im gerenderten Skript erscheint kein Logo.
\textbf{Ursache:} Pfad in \texttt{\_brand.yml} stimmt nicht, oder Datei
fehlt im \texttt{images/}-Ordner. \textbf{Lösung:} Pfad relativ zur
Projektwurzel (\texttt{images/th-koeln-logo.png}), Datei dort
tatsächlich ablegen.

\section{Wenn nichts hilft --- Tutor
fragen}\label{wenn-nichts-hilft-tutor-fragen}

\begin{tcolorbox}[enhanced jigsaw, opacitybacktitle=0.6, leftrule=.75mm, title=\textcolor{quarto-callout-tip-color}{\faLightbulb}\hspace{0.5em}{Tutor öffnen}, colframe=quarto-callout-tip-color-frame, arc=.35mm, toprule=.15mm, colbacktitle=quarto-callout-tip-color!10!white, bottomrule=.15mm, colback=white, rightrule=.15mm, breakable, toptitle=1mm, bottomtitle=1mm, opacityback=0, titlerule=0mm, coltitle=black, left=2mm]

\href{https://chatgpt.com/g/g-69fec93348108191af1ffb8ce0cf6712-tutor-interaktive-skripte}{Tutor
„Interaktive Skripte'' →}

Vorschlag-Prompt:

\begin{quote}
Ich versuche, ein Quarto-Skript zu rendern. Folgender Fehler tritt auf:

{[}Fehlermeldung kopieren{]}

Mein Setup: {[}Betriebssystem{]}, Quarto-Version {[}X.X.X{]}, R-Version
{[}X.X.X{]}, RStudio-Version {[}X.X.X{]}. Ich habe schon versucht:
{[}\ldots{]}.

Schreib mir die wahrscheinlichste Ursache in einem Satz, dann den
nächsten Schritt zum Beheben. Wenn du mehr wissen musst, frag.
\end{quote}

\end{tcolorbox}

\chapter{KI-Tutor: Link und
System-Prompt}\label{ki-tutor-link-und-system-prompt}

Der Workshop-Tutor und die Anleitung zum eigenen Tutor

\hfill\break

\subsubsection{Was dieser Anhang
leistet}\label{was-dieser-anhang-leistet-1}

Du findest hier den \textbf{fertigen Tutor} für diesen Workshop (ein
Klick) und den vollständigen \textbf{System-Prompt}, mit dem du dir
denselben Tutor in der AcademicCloud, im eigenen ChatGPT-Account oder
als statisches Widget nachbauen kannst.

\section{Der vorbereitete Workshop-Tutor (sofort
nutzbar)}\label{der-vorbereitete-workshop-tutor-sofort-nutzbar}

\begin{tcolorbox}[enhanced jigsaw, opacitybacktitle=0.6, leftrule=.75mm, title=\textcolor{quarto-callout-tip-color}{\faLightbulb}\hspace{0.5em}{Tutor öffnen}, colframe=quarto-callout-tip-color-frame, arc=.35mm, toprule=.15mm, colbacktitle=quarto-callout-tip-color!10!white, bottomrule=.15mm, colback=white, rightrule=.15mm, breakable, toptitle=1mm, bottomtitle=1mm, opacityback=0, titlerule=0mm, coltitle=black, left=2mm]

\href{https://chatgpt.com/g/g-69fec93348108191af1ffb8ce0cf6712-tutor-interaktive-skripte}{\textbf{Tutor
„Interaktive Skripte''} →}

Ein Klick --- der GPT öffnet sich in ChatGPT. Login mit kostenlosem
ChatGPT-Konto. Der System-Prompt ist eingebaut.

\end{tcolorbox}

So baust du den Link in dein eigenes Skript ein:

\begin{Shaded}
\begin{Highlighting}[]
\CommentTok{[}\OtherTok{Frag den Skript{-}Tutor →}\CommentTok{](https://chatgpt.com/g/g{-}69fec93348108191af1ffb8ce0cf6712{-}tutor{-}interaktive{-}skripte)}\NormalTok{\{target="\_blank"\}}
\end{Highlighting}
\end{Shaded}

In jedem Kapitel als Callout am Ende:

\begin{Shaded}
\begin{Highlighting}[]
\NormalTok{::: \{.callout{-}tip\}}
\FunctionTok{\#\# Tutor öffnen}
\CommentTok{[}\OtherTok{Tutor „Interaktive Skripte" →}\CommentTok{](https://chatgpt.com/g/g{-}69fec93348108191af1ffb8ce0cf6712{-}tutor{-}interaktive{-}skripte)}\NormalTok{\{target="\_blank"\}}
\NormalTok{:::}
\end{Highlighting}
\end{Shaded}

\section{Eigenen Tutor anlegen --- drei
Hosting-Wege}\label{eigenen-tutor-anlegen-drei-hosting-wege}

\subsection{Weg 1 --- eigenes Custom GPT auf
ChatGPT}\label{weg-1-eigenes-custom-gpt-auf-chatgpt}

\begin{enumerate}
\def\labelenumi{\arabic{enumi}.}
\tightlist
\item
  \url{https://chatgpt.com} öffnen, oben links \textbf{„Explore GPTs''}
  → \textbf{„+ Create''}.
\item
  \textbf{Configure}-Tab öffnen, Name vergeben, System-Prompt unten
  kopieren und einfügen.
\item
  \textbf{Save → Share} → „Anyone with the link'' oder „Public''.
\item
  Sharing-Link in dein Skript einbauen (siehe oben).
\end{enumerate}

\subsection{Weg 2 --- AcademicCloud (Hochschul-Login,
EU-Datenhaltung)}\label{weg-2-academiccloud-hochschul-login-eu-datenhaltung}

\begin{enumerate}
\def\labelenumi{\arabic{enumi}.}
\tightlist
\item
  \url{https://chat.academiccloud.de} öffnen, einloggen.
\item
  Neuen \textbf{Custom GPT} anlegen (Modell: aktuelles Top-Modell).
\item
  System-Prompt einfügen.
\item
  \textbf{Sharing aktivieren} (öffentlich oder „mit Link'').
\item
  Link einbauen.
\end{enumerate}

\subsection{Weg 3 --- statisches Widget (nur für
Fortgeschrittene)}\label{weg-3-statisches-widget-nur-fuxfcr-fortgeschrittene}

Eine HTML-Datei mit eingebautem Prompt, die direkt eine
OpenAI-kompatible API anspricht --- Vorbild:
\texttt{\_relevance-widget.qmd} in
\url{https://github.com/TH-Koln-Bartnik/genai4teaching}. Setzt
Server-Proxy oder Tagesschlüssel voraus, deshalb hier nicht im Detail.

\section{System-Prompt (Copy-Paste)}\label{system-prompt-copy-paste}

\begin{verbatim}
Du bist ein geduldiger Quarto-Coach für Hochschullehrende ohne IT-Vorkenntnisse.
Du begleitest Teilnehmende eines eintägigen Workshops, in dem sie ein
eigenes interaktives Skript mit Quarto, RStudio, GitHub Desktop und
GitHub Pages erstellen. Beispielthema ihres Skripts: eine kurze
Einführung ins wissenschaftliche Arbeiten (Wissenschaftliche Einstellung,
Recherche und Quellenqualität, oder Transparent Schreiben).

DEIN VERHALTEN
- Beginne jede Antwort mit einer Lay-Erklärung in ein bis zwei Sätzen.
  Was ist das, warum gerade jetzt? Erst danach kommt der technische Schritt.
- Frage zu Beginn: "Wo stehst du gerade?" Wenn du keinen Stand kennst,
  frage nach einer kurzen Beschreibung.
- Gib nie mehr als ZEHN Zeilen Code auf einmal. Lieber kleinere Schritte.
- Erkläre alle Pfade explizit (relativer Pfad, absoluter Pfad, was bedeutet das).
- Verwende deutsche Sprache. Bei Fachbegriffen: einmalig auf Englisch
  in Klammern, dann auf Deutsch weiter.
- Lasse Fehler nicht klein erscheinen. Wenn etwas typisch falsch läuft,
  benenne es offen und erkläre warum.

WORKSHOP-INHALT, AUF DEN DU DICH BEZIEHST
- Werkzeuge: Quarto (Version 1.5+), RStudio Desktop, GitHub Desktop,
  Zotero mit Better BibTeX, optional H5P.
- Stil: TH-Köln-Brand (rot #c81e0f, magenta #b43092, orange #ea5a00),
  scharfe Ecken, Arial.
- Struktur: Quarto-Buch mit Kapiteln in einem parts/-Ordner, output-dir: docs,
  Veröffentlichung über GitHub Pages aus dem main-Branch, Folder /docs.
- Inhaltsbeispiel: Einführung wissenschaftliches Arbeiten, drei Stränge:
  Wissenschaftliche Einstellung / Recherche und Quellenqualität /
  Transparent Schreiben (basierend auf Bergstrom & West, "Calling Bullshit").

DEINE GRENZEN
- Wechsle nicht auf andere Statische-Seiten-Generatoren (Hugo, Jekyll).
  Bleib bei Quarto.
- Antworte nicht zu allgemeinen IT-Fragen außerhalb des Skript-Workflows.
- Behaupte nichts, was du nicht sicher weißt. Bei Unsicherheit: sag das.
- Erfinde keine Quellen, keine Pakete, keine API-Endpunkte.

ZIEL
Am Ende des Workshops haben die Teilnehmenden ein eigenes, im Netz
veröffentlichtes Skript mit Inhaltsverzeichnis, Zotero-Zitaten,
TH-Köln-Stil, einem interaktiven Diagramm und einem iFrame-Element.
Hilf ihnen, dorthin zu kommen — Schritt für Schritt.
\end{verbatim}

\section{Drei Kurz-Prompts für typische
Aufgaben}\label{drei-kurz-prompts-fuxfcr-typische-aufgaben}

\subsection{A --- Fehler erklären}\label{a-fehler-erkluxe4ren}

\begin{verbatim}
Erkläre mir diesen Fehler aus dem Quarto-Render-Log in einem Satz,
dann den wahrscheinlichsten nächsten Schritt:

[Fehlermeldung einfügen]
\end{verbatim}

\subsection{B --- Code reviewen}\label{b-code-reviewen}

\begin{verbatim}
Ich habe folgenden YAML/Code in meinem Quarto-Skript. Schau, ob
Fehler oder Stolperfallen drin sind. Antworte mit drei kurzen Punkten:
was funktioniert, was nicht, was würdest du anders machen.

[Code einfügen]
\end{verbatim}

\subsection{C --- Kapitel ergänzen}\label{c-kapitel-erguxe4nzen}

\begin{verbatim}
Schlage mir für mein Kapitel zum Thema [Strang A/B/C] eine ergänzende
Übung vor — kurz, hands-on, mit konkretem Code-Snippet zum Einbauen.
Zwei bis drei Minuten Bearbeitungszeit für die Studierenden.
\end{verbatim}




\end{document}
